\documentclass[11pt]{article}
\usepackage[a4paper,margin=27mm,headheight=15pt]{geometry}
\usepackage[T1]{fontenc}
\usepackage[utf8]{inputenc}
\usepackage{lmodern,microtype}
\usepackage{amsmath,amssymb,amsthm,mathtools,mathrsfs}
\usepackage{booktabs,array,longtable,tabularx,enumitem}
\usepackage{xcolor,fancyhdr,xurl,hyperref}
\definecolor{paperink}{HTML}{20384B}
\hypersetup{colorlinks=true,linkcolor=paperink,citecolor=paperink,urlcolor=paperink,
  pdftitle={Projector-Resolved Seed Calculus IX: Consolidated Edition},
  pdfauthor={Komorebi P.},pdfsubject={Arithmetic packets, dyadic towers, polynomial lifting and analytic regularity}}
\newcommand{\N}{\mathbb N}
\newcommand{\Z}{\mathbb Z}
\newcommand{\Q}{\mathbb Q}
\newcommand{\F}{\mathbb F}
\newcommand{\Zp}{\mathbb Z_p}
\newcommand{\Qp}{\mathbb Q_p}
\DeclareMathOperator{\coker}{coker}
\DeclareMathOperator{\im}{im}
\DeclareMathOperator{\rank}{rank}
\newcommand{\lean}[1]{{\footnotesize\nolinkurl{#1}}}
\newtheorem{theorem}{Theorem}[section]
\newtheorem{proposition}[theorem]{Proposition}
\newtheorem{lemma}[theorem]{Lemma}
\newtheorem{corollary}[theorem]{Corollary}
\theoremstyle{definition}
\newtheorem{definition}[theorem]{Definition}
\newtheorem{example}[theorem]{Example}
\theoremstyle{remark}
\newtheorem{remark}[theorem]{Remark}
\numberwithin{equation}{section}
\setlength{\parskip}{3pt}
\setlength{\emergencystretch}{2em}
\setlist{itemsep=2pt,topsep=5pt}
\allowdisplaybreaks[1]
\pagestyle{fancy}
\fancyhf{}
\fancyhead[L]{\small Projector-Resolved Seed Calculus IX}
\fancyhead[R]{\small Consolidated edition}
\fancyfoot[L]{\small Komorebi P.}
\fancyfoot[R]{\thepage}

\title{\textbf{Projector-Resolved Seed Calculus IX}\\[5pt]
  \Large Arithmetic Packets, Dyadic Towers,\\ and Local Polynomial Deformation}
\author{Komorebi P.}
\date{Consolidated edition\quad 8 September 2026}

\begin{document}
\maketitle
\begin{abstract}
This paper develops arithmetic packet coordinates and their passage to local
reconstruction. It consolidates the original packet calculus, the completion
of the stated Nishibuchi congruence conjecture, and the subsequent formal
extensions. The arithmetic part distinguishes normalized and unnormalized
Euler packets, proves amplified difference bounds, gives exact dyadic
displacements and least packet periods for four residue addresses, and records
an optimal classical dyadic reconstruction tower. The lifting part identifies
finite obstruction classes, shared corrections and correction fibers, constructs
compatible polynomial towers, and proves quantitative persistence under
simultaneous coefficient and parameter deformation. An analytic implicit
function theorem for the actual polynomial system supplies local solution
graphs and their strict derivatives, with comparison to every continuous local
root selection through the same nonsingular point. The proofs and their
coefficient domains are linked to the preserved Lean sources. General singular
persistence, unrestricted presentation invariance, global descent and further
exact valuation questions remain open within this program.
\end{abstract}

\section{Scope and conventions}
\label{sec:scope}
The central objects are sequences and polynomial equations equipped with
arithmetic coordinates. Residue projectors select addresses in a sequence,
finite differences measure compatibility between nearby addresses, and
valuation bounds control the passage from finite precision to a $p$-adic
limit. These constructions provide a common language for the Euler packet
results and the later lifting results. Each theorem nevertheless retains its
own coefficient ring, topology and nonsingularity hypotheses.

This edition integrates the original arithmetic framework with the completed
extensions. The stated external Nishibuchi conjecture is proved in its full
specified parameter range. The stronger dyadic conclusions and the polynomial
deformation theorems are stated separately, with the extra conclusions and
their precise domains visible. The native conjectures retain their original
identifiers, 13.2 and 13.3, when discussed as research questions, although the
section numbering of this consolidated manuscript is new.

Throughout, $\N=\{0,1,2,\ldots\}$. When $p$ is prime, $\Zp$ denotes the
$p$-adic integers and $\Qp$ their fraction field. We normalize $v_p(p)=1$.
For a rational or $p$-adic quantity, a displayed lower valuation bound is
understood with the mathematical convention $v_p(0)=+\infty$.
The corresponding formal predicates are zero-safe: the quantity is zero or
its valuation meets the stated bound. An exact finite valuation assertion
requires a nonzero quantity, and this is stated where used. Congruence in
$\Zp$ means ideal divisibility. Thus
\[
 x\equiv y\pmod{p^r}\quad\Longleftrightarrow\quad
 x-y\in p^r\Zp.
\]
For vectors, ideal membership and congruence are coordinatewise.

Write $\Delta_u f(n)=f(n+u)-f(n)$ and use $\Delta_u^k$ for its $k$-fold
iterate. Euler coefficients indexed by ordinary degree and by packet row are
distinguished explicitly. A packet displacement by $u$ rows is a degree
displacement by $Nu$ in a packet of width $N$.

The theorem statements and mathematical arguments form the main text.
The accompanying source maps identify the relevant Lean declarations and
the exact archived source files. Verification dates and scopes are retained:
the latest 91-check polynomial build did not rebuild all earlier arithmetic
developments. The consolidation adds no new Lean theorem and makes no new
historical-priority claim. The general analytic inverse machinery used in the
last mathematical section is established Mathlib infrastructure.

\clearpage
{\small\setlength{\parskip}{0pt}\tableofcontents}
\clearpage
\section{Arithmetic packets and coefficient transport}
\label{sec:foundations}

The coefficient object in this paper is a sequence, a formal power series, or a
finite recurrence state. Residue selection changes which indices are visible.
Triangular transport changes the coordinates of the same object. A valuation
filtration then determines how many coordinates suffice at a given precision.
These operations give the common framework for the Euler congruences and for
the later finite lifting constructions.

The original arithmetic seed notation
\[
\mathcal S=(A,a_0,\mathcal O,\mathcal R,\{P_{N,j}\},w,\tau_p)
\]
records a coefficient ring or module, a seed, allowed operators, relations,
residue selectors, a depth weight, and a topology. It is useful descriptive
notation. It does not, by itself, define a category of seed presentations or a
notion of equivalence between them. The theorems below concern the concrete
objects specified in their statements. Section~\ref{sec:open} explains the
additional definitions needed for the original native conjectures.

The Lean declarations cited in this section have namespace prefix
\lean{KomorebiPaperIX}. Their exact modules and source hashes are recorded in
the accompanying statement map. They belong to the preserved baseline
formalization. The consolidation introduces no new Lean proof of these results.

\subsection{Residue selectors and exact reconstruction}

Let $R$ be a commutative ring, $a:\N\to R$, $N\geq1$, and $0\leq j<N$.
Define
\[
(P_{N,j}a)(n)=
\begin{cases}
a(n),&n\equiv j\pmod N,\\
0,&\text{otherwise}.
\end{cases}
\]
This definition keeps the original index set. The compressed subsequence
$a(Nq+j)$, used below for recurrence packets, is a different representation of
the selected coefficients.

\begin{theorem}[Residue reconstruction]
\label{thm:residue-reconstruction}
The selectors satisfy
\[
P_{N,j}^2=P_{N,j},\qquad P_{N,i}P_{N,j}=0\quad(i\ne j),
\qquad \sum_{j=0}^{N-1}P_{N,j}=I.
\]
If $N\mid M$ and $M>0$, then
\[
P_{N,j}a=\sum_{\substack{0\leq r<M\\r\equiv j\pmod N}}P_{M,r}a.
\]
The same identities hold coefficientwise for actual formal power series.
\end{theorem}

\begin{proof}
At a fixed index $n$, exactly one standard residue label equals $n\bmod N$.
Consequently a repeated selector changes nothing, two different selectors have
zero composite, and the sum retains $a(n)$ exactly once. If $N\mid M$, reducing
$n\bmod M$ modulo $N$ gives $n\bmod N$. This proves the refinement identity
at every coefficient. Equality of all coefficients proves the formal-series
versions.
\end{proof}

\begingroup\small\raggedright
\noindent\emph{Formal correspondence.}
\lean{ResiduePackets.packet_idempotent},
\lean{ResiduePackets.packet_orthogonal},
\lean{ResiduePackets.packet_reconstruction}, and
\lean{ResiduePackets.packet_refinement}; the series variants have the
prefix \lean{series_packet} in the same module.
The algebraic statements need no convergence hypothesis.
\par\endgroup

\subsection{Formal exponential coordinates and Hurwitz multiplication}

Over a characteristic-zero field $K$, put
\[
\mathcal E(a)=\sum_{n\geq0}\frac{a_n}{n!}t^n,
\qquad \mathcal D(A)_n=n![t^n]A.
\]
These are inverse maps between sequences and $K[[t]]$. We use them as
\emph{formal} exponential coordinates. A statement about their coefficients
does not assert convergence of an exponential generating function on an
analytic domain.

\begin{theorem}[Constructed unit-triangular coordinates]
\label{thm:egf-coordinates}
Let $h:\N\to K$ satisfy $h_0=1$, write $H=\mathcal E(h)$ and
$H^{-1}=\mathcal E(k)$, and define
\[
U_Ha=\mathcal D(H\mathcal E(a)).
\]
Then $U_H$ is an equivalence, with inverse
$U_H^{-1}a=\mathcal D(H^{-1}\mathcal E(a))$. Its coefficient formulas are
\[
(U_Ha)_n=\sum_{r=0}^{n}\binom nr h_r a_{n-r},\qquad
(U_H^{-1}a)_n=\sum_{r=0}^{n}\binom nr k_r a_{n-r}.
\]
The resulting triangular matrix has diagonal entries $1$. The transported
selectors
\[
\widehat P_{N,j}=U_HP_{N,j}U_H^{-1}
\]
are idempotent and pairwise orthogonal, and their sum is the identity.
\end{theorem}

\begin{proof}
The factorials are nonzero in $K$, so encoding and decoding cancel
coefficientwise. The constant coefficient of $H$ is $1$; hence its formal
power-series inverse exists. Multiplication by $H$ and $H^{-1}$ gives the two
inverse laws. The ordinary product coefficient formula, multiplied by $n!$,
gives the displayed binomial sums. In row $n$, the coefficient of $a_n$ is
$h_0=1$ and coefficients of $a_q$ with $q>n$ vanish. Conjugating the identities
of Theorem~\ref{thm:residue-reconstruction} by the constructed equivalence proves
the transported projector identities.
\end{proof}

\begingroup\small\raggedright
\noindent\emph{Formal correspondence.}
\lean{EGF.decode_encode}, \lean{EGF.encode_decode},
\lean{EGF.egf_product}, \lean{EGF.coordinateEquiv},
\lean{EGF.coordinate_forward}, \lean{EGF.coordinate_inverse},
and \lean{EGF.transported_reconstruction}, together with the idempotence and
orthogonality declarations in \lean{EGFCoordinates}.
\par\endgroup

Factorial division is inappropriate in characteristic $p$. The coefficient
operation itself remains meaningful over every commutative ring:
\[
(f\star g)_n=\sum_{r=0}^{n}\binom nr f_r g_{n-r},\qquad
(D^d f)_n=f_{n+d},\qquad \mathbf1_n=\begin{cases}1,&n=0,\\0,&n>0.\end{cases}
\]
Here $\star$ is Hurwitz multiplication. The symbol $D^d$ denotes a coefficient
shift, which agrees with formal differentiation after characteristic-zero
encoding.

\begin{theorem}[Hurwitz product rules]
\label{thm:hurwitz-product}
For $d\geq0$,
\[
D^d(f\star g)=\sum_{r=0}^{d}\binom dr(D^r f)\star(D^{d-r}g).
\]
If $R$ has prime characteristic $p$, then
\[
D^p(f\star g)=f\star D^p g+D^p f\star g.
\]
Over a field, if $g_0=1$, a sequence $b$ is constructed recursively with
$b\star g=\mathbf1$. If also $D^p g=Dg$ in characteristic $p$, then
$D^p b=Db$.
\end{theorem}

\begin{proof}
Pascal's identity proves the first product rule, and induction proves the
higher rule. The intermediate binomial coefficients $\binom pr$ vanish in
characteristic $p$. For the inverse, solve row $n$ using the coefficient of
$b_n$, which is $g_0=1$:
\[
b_n=\mathbf1_n-\sum_{r<n}\binom nr b_r g_{n-r}.
\]
Subtract the first and $p$th product rules for $b\star g=\mathbf1$.
The assumed identity for $g$ leaves
$(D^p b-Db)\star g=0$. A second triangular induction shows that convolution
with a sequence having constant coefficient $1$ has trivial kernel.
\end{proof}

\begingroup\small\raggedright
\noindent\emph{Formal correspondence.}
\lean{Hurwitz.shift_convolution}, \lean{Hurwitz.prime_shift_convolution},
\lean{Hurwitz.convolution_left_eq_zero},
\lean{HurwitzInverse.inverse_convolution}, and
\lean{HurwitzInverse.inverse_prime_difference}.
This argument uses coefficient sequences directly. It never interprets
$1/n!$ as an element of $\F_p$ when $p\mid n!$.
\par\endgroup

\subsection{Triangular inversion and a finite precision filtration}

\begin{theorem}[Forward substitution and denominator clearing]
\label{thm:triangular-foundation}
Let $T_{nq}$ be coefficients over a field with $T_{nn}\ne0$ for every $n$.
For any sequence $b$, define
\[
a_n=\frac{b_n-\sum_{q<n}T_{nq}a_q}{T_{nn}}.
\]
This constructs the unique solution of
$\sum_{q=0}^{n}T_{nq}a_q=b_n$ for all $n$. If $T_{nq}$ and $b_n$ are integers
and the solution is taken over $\Q$, then
\[
\left(\prod_{r=0}^{n}T_{rr}\right)a_n\in\Z.
\]
\end{theorem}

\begin{proof}
The recursion refers only to smaller indices. Substitution proves each row.
If a second solution agrees through row $n-1$, subtracting row $n$ and
cancelling $T_{nn}$ proves equality at $n$. For integrality, use induction and
promote every earlier clearing product to $\prod_{r<n}T_{rr}$. Multiplication
by the additional factor $T_{nn}$ cancels the only new division.
\end{proof}

\begingroup\small\raggedright
\noindent\emph{Formal correspondence.}
\lean{Triangular.solve_row}, \lean{Triangular.solve_unique}, and
\lean{Triangular.solve_denominator_integral}. This is a denominator bound,
not an assertion that the unscaled solution is integral.
\par\endgroup

To isolate the filtration argument, let $p$ be any element of a commutative
ring and let $w:\N\to\N\cup\{\infty\}$. Set $p^{[r]}=p^r$ for finite $r$
and $p^{[\infty]}=0$. The latter notation means the coefficient vanishes.

\begin{theorem}[Filtered finite reconstruction]
\label{thm:filtered-foundation}
For every $n,L\in\N$ and sequences $c,S$ in that ring,
\[
\begin{split}
&\sum_{r=0}^{n}\binom nr c_rp^{[w(r)]}S_{n-r}\\
&\hspace{8mm}\equiv
\sum_{\substack{0\leq r\leq n\\w(r)<L}}
\binom nr c_rp^{[w(r)]}S_{n-r}\pmod{p^L}.
\end{split}
\]
Here congruence means divisibility of the difference by $p^L$ in the
coefficient ring. In $\Zp$, it means membership in $p^L\Zp$.
\end{theorem}

\begin{proof}
The difference of the two sums consists exactly of terms with $w(r)\geq L$.
Every such term is divisible by $p^L$, including the zero terms with infinite
weight. Their finite sum has the same divisibility.
\end{proof}

\begingroup\small\raggedright
\noindent\emph{Formal correspondence.}
\lean{Filtered.filtered_triangular_reconstruction} and
\lean{Filtered.padic_filtered_reconstruction}.
For weights $2r$, retaining $r\leq M$ gives the bound $p^{2M+2}$ through
\lean{Filtered.even_weight_tail}. A sharper kernel weight can improve this
bound, as in the classical dyadic reconstruction later in the paper. Each
coefficient row is finite, so no hypothesis $w(r)\to\infty$ is needed here.
\par\endgroup

\subsection{Finite Mahler coordinates and weighted differences}

For $f:\N\to R$, define the step-$h$ forward difference and its iterates by
\[
\Delta_h f(n)=f(n+h)-f(n),\qquad \Delta_h^0 f=f,
\]
and put $c_r(f)=\Delta_1^r f(0)$. These are actual differences of $f$, not
freely supplied coefficients.

\begin{theorem}[Finite Newton--Mahler identities]
\label{thm:mahler-foundation}
For every $n,s,h\in\N$,
\[
\begin{aligned}
f(n)&=\sum_{r=0}^{n}c_r(f)\binom nr,\\
c_r(\Delta_1^s f)&=c_{s+r}(f),\\
\Delta_1^s f(n)&=\sum_{j=0}^{n}c_{s+j}(f)\binom nj,\\
\Delta_h^s f(n)&=\sum_{j=0}^{s}(-1)^{s-j}\binom sj f(n+jh).
\end{aligned}
\]
The finite interpolant
$I_Nf(n)=\sum_{r=0}^{N}c_r(f)\binom nr$ agrees with $f$ for $n\leq N$.
It has the same order-$s$ unit-step difference at $n$ whenever $n+s\leq N$.
For arbitrary coefficients $c_0,\ldots,c_N$, the corresponding binomial
interpolant has exactly those coefficients and zero coefficients above $N$.
\end{theorem}

\begin{proof}
The shift operator is $I+\Delta_1$. Expanding its $n$th power and evaluating
at zero gives reconstruction. Composing difference operators gives the
coefficient shift and applying reconstruction to $\Delta_1^s f$ gives the
third identity. The binomial expansion of $(E_h-I)^s$ gives the explicit
stencil. Binomial coefficients vanish when their lower index exceeds $n$,
which proves the interpolation claim. Equality of every point of the stencil
proves equality of differences. Repeated finite differences of $\binom nr$
at zero give the exact coefficient extraction.
\end{proof}

\begingroup\small\raggedright
\noindent\emph{Formal correspondence.}
\lean{Mahler.finite_mahler_reconstruction},
\lean{Mahler.coefficient_difference},
\lean{Mahler.difference_from_coefficients},
\lean{Mahler.difference_explicit},
\lean{Mahler.finite_sample_difference}, and
\lean{Mahler.coefficient_interpolant}.
Finite samples determine only stencils contained in the sample interval.
They do not establish a universal tail estimate.
\par\endgroup

For $x\in\Zp$, use the zero-safe valuation
\[
\operatorname{ord}_p(x)=
\begin{cases}v_p(x),&x\ne0,\\\infty,&x=0.\end{cases}
\]
Thus $L\leq\operatorname{ord}_p(x)$ is equivalent to $x\in p^L\Zp$.

\begin{theorem}[Weighted difference bound]
\label{thm:weighted-mahler-foundation}
Let $f:\N\to\Zp$ and let $w:\N\to\N\cup\{\infty\}$ satisfy
$w(s+j)\leq\operatorname{ord}_p(c_{s+j}(f))$ for $0\leq j\leq n$.
Then
\[
\operatorname{ord}_p(\Delta_1^s f(n))\geq
\min_{0\leq j\leq n}
\left\{w(s+j)+v_p\binom nj\right\}.
\]
\end{theorem}

\begin{proof}
Apply the nonarchimedean inequality to the finite coefficient formula in
Theorem~\ref{thm:mahler-foundation}. For $0\leq j\leq n$, the integer
$\binom nj$ is positive. The order of its product with $c_{s+j}(f)$ is the
sum of the two orders, including the case where that coefficient is zero.
\end{proof}

\begingroup\small\raggedright
\noindent\emph{Formal correspondence.}
\lean{PadicBounds.nat_le_order_iff},
\lean{PadicBounds.weighted_difference_bound}, and
\lean{PadicBounds.weighted_difference_divisibility}.
These are finite algebraic identities and bounds. An expansion on all of
$\Zp$ additionally requires an appropriate convergence theorem; no such
infinite expansion is inferred from finite interpolation alone.
\par\endgroup

\subsection{Unit modes and the dyadic distinction}

\begin{theorem}[Differences of unit modes]
\label{thm:unit-modes-foundation}
Over every commutative ring,
\[
\Delta_h^s(u^n)=u^n(u^h-1)^s.
\]
For any prime $p$, units $u_i\in\Zp^\times$, coefficients $a_i\in\Zp$, and
a finite index set $I$, the sequence $f(n)=\sum_{i\in I}a_i u_i^n$ satisfies
\[
\Delta_{p-1}^s f(n)\in p^s\Zp
\qquad(n,s\in\N).
\]
\end{theorem}

\begin{proof}
Induction gives the geometric identity. Reduction of a unit to $\F_p^\times$
gives $u_i^{p-1}=1$ modulo $p$. Taking $s$ powers supplies $p^s$ in each
summand. The actual difference operator is linear over finite sums.
\end{proof}

\begingroup\small\raggedright
\noindent\emph{Formal correspondence.}
\lean{UnitModes.geometric_difference},
\lean{UnitModes.unit_fermat_divisibility}, and
\lean{UnitModes.finite_unit_modes_divisibility}.
The result includes $p=2$, where the displayed step is $1$. It makes no
claim about arbitrary infinite sums of modes.
\par\endgroup

\begin{theorem}[Finite unit coordinates]
\label{thm:finite-unit-foundation}
Let $p$ be prime, $L\geq0$, and
$a\in(\Z/p^{L+1}\Z)^\times$. Define
\[
\omega_L(a)=a^{p^L},\qquad \langle a\rangle_L=a\omega_L(a)^{-1}.
\]
Then
\[
a=\omega_L(a)\langle a\rangle_L,\qquad
\omega_L(a)^{p-1}=1,\qquad \langle a\rangle_L\equiv1\pmod p.
\]
The two factors are uniquely determined by these properties and are compatible
with reduction to lower precision.
For $a\in(\Z/2^{L+2}\Z)^\times$, there is instead a unique decomposition
\[
a=\epsilon(a)\langle a\rangle_4,
\qquad \epsilon(a)\in\{1,-1\},\qquad
\langle a\rangle_4\equiv1\pmod4.
\]
\end{theorem}

\begin{proof}
The unit-group exponent $p^L(p-1)$ gives the root identity for $\omega_L$.
Its residue is that of $a$, so the remaining factor has residue $1$.
If two $(p-1)$st roots have equal residue, raising them to $p^L$ makes them
equal modulo $p^{L+1}$, while each root is fixed by this exponent. This proves
uniqueness and then compatibility under reduction. In the dyadic construction,
an odd residue is exactly $1$ or $-1$ modulo $4$. Choose that sign and multiply
$a$ by it to obtain the principal factor. The two signs remain distinct
modulo $4$, which proves uniqueness.
\end{proof}

\begingroup\small\raggedright
\noindent\emph{Formal correspondence.}
\lean{FiniteUnits.finite_unit_decomposition},
\lean{FiniteUnits.decomposition_unique},
\lean{UnitRefinement.omega_compatible},
\lean{UnitRefinement.principal_compatible}, and
\lean{DyadicUnits.dyadic_unit_decomposition} with
\lean{DyadicUnits.decomposition_unique}.
The finite $p=2$ version of the first construction has trivial root-of-unity
factor. The separate modulo-$4$ construction retains the sign that reduction
modulo $2$ cannot detect.
\par\endgroup

\subsection{Recurrence packets and finite arithmetic synchronization}

\begin{theorem}[Companion jumps and packet recurrences]
\label{thm:companion-foundation}
Let $C$ be a $d\times d$ matrix over a commutative ring and
$x_n=C^n x_0$. For $N,j,k\in\N$, put $b_k=x_{Nk+j}$. Then
\[
b_k=C^j(C^N)^kx_0,\qquad b_{k+1}=C^Nb_k.
\]
If $\chi_{C^N}(X)=X^d+\sum_{r<d}q_rX^r$, then
\[
b_{k+d}=-\sum_{r=0}^{d-1}q_r b_{k+r}.
\]
Every coordinate and every linear scalar observable satisfies the same
recurrence. No invertibility or diagonalization hypothesis is required.
\end{theorem}

\begin{proof}
The jump identity is the power law for one matrix. Apply the
Cayley--Hamilton identity $\chi_{C^N}(C^N)=0$ to the actual state $b_k$.
The leading coefficient is $1$, so the final state can be isolated. A linear
observable preserves the resulting finite sum.
\end{proof}

\begingroup\small\raggedright
\noindent\emph{Formal correspondence.}
\lean{Companion.companion_jump},
\lean{Companion.packet_charpoly_recurrence}, and
\lean{Companion.packet_observable_recurrence}.
The source-sequence versions require and use the actual relation
$x_{n+1}=Cx_n$, rather than assuming the desired packet recurrence.
For the Fibonacci matrix $C=\left(\begin{smallmatrix}1&1\\1&0\end{smallmatrix}\right)$,
its jump trace is the Lucas number $L_N$ and its determinant is $(-1)^N$.
Thus
\[
F_{N(k+2)+j}=L_NF_{N(k+1)+j}-(-1)^NF_{Nk+j}.
\]
This specialization is
\lean{FibonacciPackets.fibonacci_packet_recurrence}; it is derived from the
actual Fibonacci state.
\par\endgroup

\begin{theorem}[Dyadic--odd packet factorization]
\label{thm:crt-foundation}
Let $a\geq0$ and let $M$ be a positive odd integer. The Chinese remainder map
is a ring isomorphism
\[
\Z/(2^aM)\Z\simeq\Z/2^a\Z\times\Z/M\Z.
\]
For $0\leq r<2^aM$ and any coefficient sequence $f$,
\[
P_{2^aM,r}f
=P_{2^a,r\bmod 2^a}\,P_{M,r\bmod M}f.
\]
The same identity holds for formal-series coefficient selectors.
\end{theorem}

\begin{proof}
The two moduli are coprime. An index has the prescribed combined residue
exactly when it has both component residues. Applying this equivalence to
the indicator defining each selector proves the formula.
\end{proof}

\begingroup\small\raggedright
\noindent\emph{Formal correspondence.}
\lean{CRTPackets.addressEquiv}, \lean{CRTPackets.residue_equivalence},
\lean{CRTPackets.packet_factorization}, and
\lean{CRTPackets.series_packet_factorization}.
This theorem synchronizes finite arithmetic addresses. It does not identify a
global integer object with an arbitrary family of $p$-adic objects.
\par\endgroup

\subsection{Polyphase notation retained from the recovered manuscript}
\label{subsec:polyphase-exposition}

The recovered Annals text develops a broader polyphase vocabulary. The
following elementary identities explain that vocabulary; they are exposition,
not additional Lean-certified assertions in this consolidation.
For a sequence $a$ and $M\geq1$, let $a_r^{[M]}(q)=a_{Mq+r}$ for $0\leq r<M$.
If $A(z)=\sum_n a_nz^n$ and $A_r(w)=\sum_q a_{Mq+r}w^q$ are ordinary generating
functions, then unique quotient and remainder give
\[
A(z)=\sum_{r=0}^{M-1}z^rA_r(z^M),\qquad
\bigl(a_r^{[M]}\bigr)_s^{[N]}=a_{r+Ms}^{[MN]}.
\]
For ordinary Cauchy multiplication, the phase-$r$ product has a term without a
carry and a term with one carry:
\[
C_r(w)=\sum_{u+v=r}A_u(w)B_v(w)
+w\sum_{u+v=r+M}A_u(w)B_v(w),\qquad 0\leq u,v<M.
\]
This follows by adding the two phase addresses. It is an ordinary-series
formula and must not be substituted for Hurwitz convolution.

The affine Mahler kernel is
\[
\Theta_{M,s}(r,k)=[z^r](1+z)^s\bigl((1+z)^M-1\bigr)^k
=\sum_{j=0}^{k}(-1)^{k-j}\binom kj\binom{Mj+s}{r}.
\]
It vanishes unless $k\leq r\leq Mk+s$. For a finite binomial expansion
$f(n)=\sum_{r=0}^{D}c_r\binom nr$, the $k$th Newton coefficient of
$g(q)=f(Mq+s)$ is therefore
\[
c_k(g)=\sum_{r=k}^{\min(D,Mk+s)}c_r\Theta_{M,s}(r,k).
\]
Indeed, insert the finite stencil for $\Delta_1^k g(0)$ and interchange its
two finite sums. Coefficient extraction of the binomial expansion of
$((1+z)^M-1)^k$ gives the kernel. In the binary zero-address case,
\[
\Theta_{2,0}(k+j,k)=2^{k-j}\binom kj\qquad(0\leq j\leq k),
\]
since $(2z+z^2)^k=z^k(2+z)^k$.
The general exact valuation problem is still distinct from these identities.
The recovered infinite analytic extensions and stronger historical identities
are accounted for in Section~\ref{subsec:historical-scope}; this finite
exposition does not certify those extensions.

\section{Congruential Euler packets and the completed arithmetic theorems}
\label{sec:arithmetic}

The arithmetic results concern two related but distinct constructions.  The
congruential Euler family is indexed by a modulus $N$ and an address $j$.
Its exceptional dyadic member has $N=4$ and four addresses.  The classical
signed Euler numbers arise from $N=2$, $j=0$.  We retain these indices
throughout so that the exact four-packet periods and the classical edge
truncations cannot be confused.

\subsection{Original coefficients, normalized packets, and precision}

\begin{definition}[Congruential Euler coefficients]\label{def:arith-euler}
For $N\ge1$ and $0\le j<N$, define rational coefficients by the formal identity
\begin{equation}\label{eq:arith-euler-egf}
 \sum_{d\ge0}E_d^{(N,j)}\frac{t^d}{d!}
 =\left(\sum_{n\ge0}\frac{t^{Nn}}{(Nn+j)!}\right)^{-1}.
\end{equation}
The reciprocal exists because its denominator has constant coefficient
$1/j!\ne0$.  Set
\[
 U_{N,j}(n)=E_{Nn}^{(N,j)},\qquad C_{N,j}(n)=\frac{U_{N,j}(n)}{j!}.
\]
Thus $C_{N,j}(0)=1$ and $U_{N,j}(0)=j!$.  All coefficients outside the
multiples of $N$ vanish.
\end{definition}

The normalized source and its inverse are
\begin{equation}\label{eq:arith-kernel}
 K_{N,j}(t)=j!\sum_{n\ge0}\frac{t^{Nn}}{(Nn+j)!},\qquad
 \mathcal P_{N,j}(t)=\sum_{n\ge0}C_{N,j}(n)\frac{t^{Nn}}{(Nn)!}.
\end{equation}
Their product is one.  Clearing factorials in its coefficient of $t^{Nn}$
gives the exact triangular system
\begin{equation}\label{eq:arith-triangular}
 \sum_{q=0}^{n}\binom{Nn+j}{Nq}C_{N,j}(q)=\delta_{n0}.
\end{equation}
In particular, for $n>0$,
\[
 C_{N,j}(n)=-\binom{Nn+j}{Nn}^{-1}
       \sum_{q=0}^{n-1}\binom{Nn+j}{Nq}C_{N,j}(q).
\]
The diagonal binomial coefficient is positive, so this recurrence determines
the sequence uniquely.  In the formal development the recurrence is the
initial construction.  The reciprocal identity is subsequently proved from
it, as is the identification with the literal degree-indexed coefficients.
Consequently computations using the recurrence concern the coefficients of
\eqref{eq:arith-euler-egf} themselves.

The relevant declarations are
\lean{KomorebiPaperIX.Congruential.triangular_inverse},
\lean{KomorebiPaperIX.Congruential.packet_is_actual_inverse}, and
\lean{KomorebiPaperIX.NishibuchiCompletion.generalizedEuler_at_packet}.
The cleared identity with the original positive-address exponential section
is \lean{KomorebiPaperIX.NishibuchiCompletion.original_generating_identity}.

\begin{definition}[Differences and valuation conventions]\label{def:arith-precision}
For a sequence $f$ and $h,k,n\in\N$, put
\[
 \Delta_h f(n)=f(n+h)-f(n),\qquad
 \Delta_h^0 f(n)=f(n),\qquad
 \Delta_h^{k+1}f(n)=\Delta_h^k f(n+h)-\Delta_h^k f(n).
\]
We use the extended valuation $v_p(0)=+\infty$ in mathematical statements.
Thus $v_p(x)\ge L$, including negative integers $L$, means
$x=0$ or the ordinary valuation of the nonzero rational $x$ is at least $L$.
For integral coefficients and $L\ge0$, it is equivalent to
$x\in p^L\Z_p$.  In expressions involving a step multiplier we write
$\nu_p(u)=v_p(u)$ for $u>0$ and set $\nu_p(0)=0$.  This finite convention
makes the displayed exponents meaningful at step zero, where all
positive-order differences vanish.
\end{definition}

The explicit alternative at zero is the predicate
\lean{KomorebiPaperIX.GlobalEulerBoundary.precision}.  All divisibility
statements below are in $\Z_p$.  They do not use field divisibility in
$\Q_p$ as a substitute for a precision statement.

\subsection{The full Nishibuchi statement and its starting bound}

\begin{theorem}[Completed Nishibuchi congruence]\label{thm:arith-nishibuchi}
Let $p$ be prime, let $m>0$, and suppose that $m\mid p-1$ or $m=2$.
Put $N=mp$ and $q=\operatorname{lcm}(2,p-1)$.  For every $0\le j<N$,
all $U_{N,j}(n)$ belong to $\Z_p$.  For $r\ge1$ and $n\ge r$,
\begin{equation}\label{eq:arith-nishibuchi}
 E_{Nn+qp^r}^{(N,j)}\equiv E_{Nn}^{(N,j)}\pmod{p^r}.
\end{equation}
Thus the positive sufficient starting index $n_0=r$ works uniformly.
If $m\mid p-1$, the stronger sufficient condition $r\le mn$ is valid.
\end{theorem}

This is the full arithmetic parameter range and displacement of Conjecture
3.9 in the fixed version of Nishibuchi's paper used for the statement
comparison.  Its formal capstone is
\lean{KomorebiPaperIX.NishibuchiCompletion.conjecture_3_9}.
The stronger divisor-branch threshold is
\lean{KomorebiPaperIX.NishibuchiCompletion.divisor_period}.
The theorem asserts a sufficient starting index.  It does not assert the
earliest index for every tuple of parameters.

We give the ingredients before assembling the proof.  The first ingredient
is the complete all-address theorem on the divisor branch.

\begin{theorem}[All-address differences on the divisor branch]
\label{thm:arith-divisor-base}
Let $p$ be prime, $m>0$, $m\mid p-1$, $N=mp$, and $0\le j<N$.
Write $h=(p-1)/m$.  For every $u,k,n\in\N$,
\begin{align}
 U_{N,j}(n)&\in\Z_p,\label{eq:arith-divisor-integral}\\
 \Delta_{uh}^{k}U_{N,j}(n)&\in p^{\min(k,mn)}\Z_p,\label{eq:arith-divisor-bound}\\
 v_p\bigl(\Delta_{uh}^{k}C_{N,j}(n)\bigr)
   &\ge\min(k,mn)-v_p(j!).\label{eq:arith-divisor-normalized}
\end{align}
The factorial correction is part of the normalized statement, including
addresses $j\ge p$.
\end{theorem}

\begin{proof}
The proof works with actual formal series in the Hurwitz integral ring
\[
 \mathcal H_p=\left\{F\in\Q_p[[t]]:
            d![t^d]F\in\Z_p\text{ for every }d\right\}.
\]
For $D=d/dt$, let $B_p=(D^{p^2}-D^p)/p$.  Call $F$ flat if
$B_p^sF\in\mathcal H_p$ for every $s\ge0$.
The proof establishes flatness of the actual series $j!\mathcal P_{N,j}$ for every
address, including the addresses at which the usual unit-pivot argument
fails.  Here is how those addresses enter.

Set
\[
 H_c(t)=\sum_{n\ge0}\frac{t^{Nn+c}}{(Nn+c)!},\qquad
 A_c(t)=\sum_{n\ge0}\frac{t^{Nn}}{(Nn+c)!},\qquad
 R_c(t)=\frac{A_{c-1}(t)}{A_c(t)}\quad(1\le c<N).
\]
The denominator of this last quotient has nonzero constant coefficient.
The exact section identities give
\[
 R_cH_c=tH_{c-1},\qquad R_c(0)=c,\qquad
 j!\mathcal P_{N,j}=\mathcal P_{N,0}\prod_{c=1}^{j}R_c.
\]
The product identity follows from cancellation of consecutive source
sections, so the desired family is not replaced by an auxiliary one.

At the singular addresses $c=ap$, $1\le a<m$, use the weighted ring
\[
 \mathcal W_p=\left\{\sum_{d\ge0}a_dT^d:
          p^{\lfloor d/p\rfloor}a_d\in\Z_p\right\}.
\]
The source proves $R_{ap}(\log(1+T))\in\mathcal W_p$.  Its integral
arithmetic input can be seen in the cyclic group ring $\Z[\Z/N\Z]$.
For its canonical element $x$ and $I=(p,x^{p-1}-1)$, the relation
$(x^{p-1})^p=1$ implies $I^p\subseteq(p)I$.  Induction yields
$I^q\subseteq(p^{v_p(q!)})$.  Grouping the falling factorial
$(x)_d=\prod_{i=0}^{d-1}(x-i)$ into consecutive blocks of length $p$
then gives an integral residue-coordinate factor
$p^{v_p(\lfloor d/p\rfloor!)}$.  The identity
\[
 v_p(d!)=\lfloor d/p\rfloor+v_p(\lfloor d/p\rfloor!)
\]
accounts for precisely the weight in $\mathcal W_p$ after dividing by $d!$.
Comparison of the section and group-ring recurrences identifies these
coordinates with the actual logarithmically substituted section.

To pass from sections to the singular ratio, remove its forced zero:
for $c=ap$, the section has the form
\[
 H_c(\log(1+T))=\frac{T^c}{c!}\,V(T),\qquad V(0)=1,
 \qquad V\in\mathcal W_p.
\]
The weight of $c!$ compensates for the coefficient shift by $c$.
The weighted ring is closed under the actual unit-series inverse and under
$L_T=\log(1+T)(1+T)d/dT$.  With
$\ell(T)=\log(1+T)/T$, interpreted by coefficient shifting, differentiation
of the displayed factorization gives
\[
 R_c(\log(1+T))=c(1+T)\ell(T)+(L_TV)V^{-1}.
\]
Since $p\ell(T)\in\mathcal W_p$ and $p\mid c$, both terms are weighted.
This is the singular estimate
\lean{KomorebiPaperIX.SingularRatioAmice.paper_singular_amice_estimate}.

The bridge from that weighted estimate to flatness uses
$\Theta=(1+T)d/dT$.  Both $\Theta$ and
$(\Theta^{p^2}-\Theta^p)/p$ preserve $\mathcal W_p$.
For the first operator, its coefficient formula is
$(d+1)a_{d+1}+da_d$; whenever the floor weight rises, the factor $d+1$
provides the missing $p$.  For the second operator, the proof uses its
full Leibniz correction and the scalar congruence
$p^2\mid b^{p^2}-b^p$.  Weighted monomials factor as
\[
 a_dT^d=(p^{\lfloor d/p\rfloor}a_d)
          (T^p/p)^{\lfloor d/p\rfloor}T^{d\bmod p}.
\]
The operator preserves these monomials and hence finite weighted
polynomials.  Each fixed output coefficient depends on only finitely many
input coefficients, which extends the result to the formal series.
Substitution $T=e^t-1$ intertwines $D$ with $\Theta$ and maps
$\mathcal W_p$ into $\mathcal H_p$, because
\[
 d![t^d]F(e^t-1)=[T^0]\Theta^dF.
\]
It follows that the actual singular ratios are flat.  This transfer is
\lean{KomorebiPaperIX.WeightedDworkTransfer.singular_ratio_flat}.

At a nonsingular address $2\le c<N$ with $p\nmid c$, differentiation of the section
identities gives
\[
 tR_c'=R_c(1+R_{c-1}-R_c).
\]
Writing $q=R_{c-1}-(c-1)$ and $y=R_c-c$, this becomes
\[
 (tD+c)y=cq+yq-y^2.
\]
On degrees divisible by $p$, the inverse of $tD+c$ divides the degree-$d$
coefficient by the genuine unit $d+c$.  Starting with $y_0=0$, form
\[
 y_{s+1}=(tD+c)^{-1}(cq+y_sq-y_s^2).
\]
Flatness and support are preserved.  Since the difference of two right
sides factors as $(z-y)(q-z-y)$ and its second factor has zero constant
term, the agreement with the actual solution gains one coefficient at
each iteration.  Consequently $[t^d]y_s=[t^d]y$ for $d\le s$.
Every coefficient of $B_p^r y$ can therefore be checked in a sufficiently
long finite iterate.  No analytic convergence hypothesis is needed.
This is the content of
\lean{KomorebiPaperIX.RiccatiFlatPropagation.riccati_solution_flat}.

The standard window supplies the initial ratio $R_1$.  Singular weighted transfer
handles $p\mid c$, and Riccati propagation handles $p\nmid c$.
Flatness is closed under products, so the exact factorization establishes
flatness of $j!\mathcal P_{N,j}$.  Extracting its supported Hurwitz coefficients
and applying the proved saturation and step-multiple lemmas yields
\eqref{eq:arith-divisor-integral}--\eqref{eq:arith-divisor-bound}.
Finally $\Delta_{uh}^kU_{N,j}=j!\Delta_{uh}^kC_{N,j}$ gives
\eqref{eq:arith-divisor-normalized}, with the zero case retained.
The assembled declarations are
\lean{KomorebiPaperIX.AllResidueEuler.all_residue_flat},
\lean{KomorebiPaperIX.AllResidueEuler.global_unnormalized_bound}, and
\lean{KomorebiPaperIX.AllResidueEuler.global_normalized_precision}.
\end{proof}

\begin{lemma}[Prime-power period from finite differences]
\label{lem:arith-newton-period}
Let $R$ be a commutative ring, $p$ a prime, $f:\N\to R$, $h,n\in\N$,
and $r\ge1$.  If
\[
 p^{\min(k,r)}\mid\Delta_h^k f(n)\quad\text{for every }k\in\N,
\]
then $p^r\mid f(n+p^{r-1}h)-f(n)$.
\end{lemma}

\begin{proof}
The finite Newton expansion gives
\[
 f(n+p^{r-1}h)-f(n)
 =\sum_{k=1}^{p^{r-1}}\binom{p^{r-1}}k\Delta_h^k f(n).
\]
Terms with $k\ge r$ have the required factor from the difference.
For $1\le k<r$, the identity
\[
 k\binom{p^{r-1}}k=p^{r-1}\binom{p^{r-1}-1}{k-1}
\]
and $v_p(k)<k$ imply
$p^r\mid\binom{p^{r-1}}k p^k$ in $\Z$.
Mapping that integer divisibility to $R$ proves the claim term by term.
The argument includes $r=1$.
The declaration is
\lean{KomorebiPaperIX.PeriodFromDifferences.prime_power_period}.
\end{proof}

\begin{proof}[Proof of Theorem~\ref{thm:arith-nishibuchi}]
On the divisor branch, $m\mid q$ and the packet step $q/m$ is an integer
multiple of $h=(p-1)/m$.  If $r\le mn$, Theorem~\ref{thm:arith-divisor-base}
supplies the hypotheses of Lemma~\ref{lem:arith-newton-period}.
Its packet displacement is $p^{r-1}q/m$.  Conversion to the original
degree is exact:
\[
 mp\left(n+p^{r-1}\frac qm\right)=mpn+qp^r.
\]
The choice $n\ge r$ implies $mn\ge r$ because $m\ge1$.

For an odd prime, $2\mid p-1$, so the option $m=2$ introduces no further
case.  The sole remaining branch is $p=2,m=2,N=4$.  The all-row theorem
proved next gives $2^k\mid\Delta_1^kC_{4,j}(n)$ at every row and address.
After multiplying by $j!$, Lemma~\ref{lem:arith-newton-period} gives packet
period $2^{r-1}$ and degree displacement $4\cdot2^{r-1}=2^{r+1}=qp^r$.
It also gives integrality.  Hence the complete statement follows.
The formal degree conversion is
\lean{KomorebiPaperIX.NishibuchiCompletion.original_degree_shift};
the all-parameter degree theorem is
\lean{KomorebiPaperIX.NishibuchiCompletion.full_period_degree}.
\end{proof}

\subsection{All-row dyadic differences and displacement amplification}

For the remainder of the four-packet discussion abbreviate
$C_j=C_{4,j}$ and $U_j=U_{4,j}$.  Removing the original row cap requires
an unshifted coefficient operator.

\begin{theorem}[All-row dyadic higher differences]\label{thm:arith-dyadic-differences}
For $0\le j<4$ and all $u,k,n\in\N$,
\begin{align}
 C_j(n)&\in\Z_2,\\
 \Delta_u^kC_j(n)&\in 2^{k(1+\nu_2(u))}\Z_2,\label{eq:arith-allrow-c}\\
 \Delta_u^kU_j(n)&\in
 2^{k(1+\nu_2(u))+v_2(j!)}\Z_2.\label{eq:arith-allrow-u}
\end{align}
There is no restriction $k\le n$ and no minimum with the row $n$.
\end{theorem}

\begin{proof}
The Hurwitz coefficients of the normalized source $K_{4,j}$ at degrees
$4n$ are
\[
 g_j(n)=\binom{4n+j}{4n}^{-1}.
\]
Their literal expressions are
\[
 \begin{aligned}
 g_0(n)&=1,&
 g_1(n)&=\frac1{4n+1},\\
 g_2(n)&=\frac1{(4n+1)(2n+1)},&
 g_3(n)&=\frac3{(4n+1)(2n+1)(4n+3)}.
 \end{aligned}
\]
Every denominator is odd.  An affine sequence with even slope satisfies
$2^k\mid\Delta_1^k f(n)$ for every $k,n$.  The exact product identity
\[
 \Delta_1^s(fg)(n)=\sum_{i=0}^s\binom si
          \Delta_1^if(n)\Delta_1^{s-i}g(n+i)
\]
proves closure of this property under products.  Applying the formula
with the factors interchanged to $fg=1$ isolates $f(n+s)\Delta_1^sg(n)$; multiplication by the actual inverse
$f(n+s)^{-1}$ and induction prove closure under pointwise unit inverses.
Thus all four $g_j$ have the property.  The definitions and inverse
identities are checked in
\lean{KomorebiPaperIX.DyadicFourKernel.source_coe} and
\lean{KomorebiPaperIX.DyadicFourKernel.source_kummer}.

For an integer charge $c$, put
\[
 A_c=\frac{D^4-c}{2}.
\]
A series is flat of charge $c$ when all $A_c^s$ of it belong to
$\mathcal H_2$.  The required product formula is
\begin{equation}\label{eq:arith-charged-product}
 A_{c+d}(FG)=(A_cF)G+F(A_dG)
       +2F'G'''+3F''G''+2F'''G'.
\end{equation}
All middle coefficients are integral.  Derivatives preserve
$\mathcal H_2$ and commute with $A_c$.  Induction on the number of
iterations in \eqref{eq:arith-charged-product} proves that multiplication
adds charges.  If $FG=1$, the same formula yields
\[
 A_{-c}G=-G^2A_cF-G(2F'G'''+3F''G''+2F'''G').
\]
Induction, starting from the independently established integrality of
the actual reciprocal $G$, proves that inversion negates the charge.
Finally $A_{c+2}=A_c-1$, so increasing a charge by two preserves flatness.
These are
\lean{KomorebiPaperIX.ChargedDworkClosure.B_mul},
\lean{KomorebiPaperIX.ChargedDworkClosure.flat_inverse}, and
\lean{KomorebiPaperIX.ChargedDworkClosure.flat_charge_shift}.

Packing $g_j$ into degrees divisible by four gives $K_{4,j}$.
Its source property is precisely charge-one flatness.  Inversion gives
charge $-1$, and the charge shift restores charge one for the actual
reciprocal $\mathcal P_{4,j}$.  If $\operatorname{dec}(F)(d)=d![t^d]F$, the exact
coefficient identity is
\[
 \operatorname{dec}(A_1^sF)(d)
       =2^{-s}\Delta_4^s\operatorname{dec}(F)(d).
\]
At $d=4n$, packing turns the right side into
$2^{-s}\Delta_1^sC_j(n)$.  This extraction requires no negative degree
and therefore holds at every $n$, including zero.  It proves the unit-step
bound \lean{KomorebiPaperIX.DyadicFourAllRows.base_step_bound}.

To incorporate $v_2(u)$, expand a displacement using finite Newton
coefficients.  For any prime $p$, $u>0$, and $1\le a\le u$,
\[
 p^{v_p(u)+1}\mid\binom ua p^a,
\]
as follows from $a\binom ua=u\binom{u-1}{a-1}$ and $v_p(a)<a$.
Repeated displacement, retaining the remaining base-difference order
at each induction step, gives a factor $p^{k(1+v_p(u))}$ from $k$
displacements.  The source supplies this induction as
\lean{KomorebiPaperIX.StepAmplification.repeated_precision}.
Applying it with $p=2$ proves \eqref{eq:arith-allrow-c}; multiplication
by $j!$ proves \eqref{eq:arith-allrow-u}.  Zero step and zero order follow
from their definitions.  The final statements are
\lean{KomorebiPaperIX.DyadicFourFinalStrengthening.normalized_step_bound}
and
\lean{KomorebiPaperIX.DyadicFourFinalStrengthening.unnormalized_step_bound}.
\end{proof}

\begin{corollary}[Step-sensitive divisor-branch bound]\label{cor:arith-amplified}
Under the hypotheses of Theorem~\ref{thm:arith-divisor-base}, for all
$u,k,n\in\N$,
\[
 \Delta_{u(p-1)/m}^{k}U_{mp,j}(n)
       \in p^{\min(k(1+\nu_p(u)),mn)}\Z_p.
\]
For $C_{mp,j}$ the corresponding valuation lower bound is
$\min(k(1+\nu_p(u)),mn)-v_p(j!)$.
\end{corollary}

\begin{proof}
Apply the same repeated-displacement theorem with the fixed cap $mn$ to
\eqref{eq:arith-divisor-bound}.  The argument preserves this cap and
improves the contribution from the order of the difference.
Then divide by the fixed factorial in the zero-safe rational formulation.
These are
\lean{KomorebiPaperIX.AllResidueStepAmplification.unnormalized_precision}
and
\lean{KomorebiPaperIX.AllResidueStepAmplification.normalized_precision}.
\end{proof}

\begin{corollary}[Dyadic periods from row zero]\label{cor:arith-allrow-period}
For $0\le j<4$, $r\ge1$, and $n\ge0$,
\[
 U_j(n+2^{r-1})\equiv U_j(n)\pmod{2^r}.
\]
For $j=2,3$ and $r\ge2$, $2^{r-2}$ is also a sufficient packet period.
The respective original degree displacements are $2^{r+1}$ and $2^r$.
\end{corollary}

\begin{proof}
Use Lemma~\ref{lem:arith-newton-period} and the unit-step case of
Theorem~\ref{thm:arith-dyadic-differences}.  At $j=2,3$ the factorial
contributes an additional power of two.  The actual degree versions are
\lean{KomorebiPaperIX.DyadicFourFinalStrengthening.all_rows_period_degree}
and
\lean{KomorebiPaperIX.DyadicFourFinalStrengthening.all_rows_half_period_degree}.
The next theorem sharpens these sufficient periods to a complete
classification.
\end{proof}

\subsection{Exact displacement valuations and least packet periods}

\begin{lemma}[Finite Newton dominance]\label{lem:arith-dominance}
Let $f:\N\to\Q$ and write $a_k=\Delta_1^kf(0)$.  Suppose that $c\in\N$,
$a_1\ne0$, $v_2(a_1)=c$, and
\[
 v_2(a_k)\ge c+1+\lfloor\log_2k\rfloor\qquad(k\ge2).
\]
Then for every $n\ge0$ and $T>0$,
\[
 f(n+T)-f(n)\ne0,\qquad
 v_2\bigl(f(n+T)-f(n)\bigr)=v_2(T)+c.
\]
\end{lemma}

\begin{proof}
Finite Newton reconstruction at an arbitrary row gives
\[
 \Delta_1^kf(n)=\sum_{i=0}^{n}\binom ni a_{k+i}.
\]
The monotonicity of $\lfloor\log_2k\rfloor$ therefore preserves the
higher-coefficient bound after translation.  In addition,
$\Delta_1f(n)-a_1\in2^{c+1}\Z_2$, so the first difference is nonzero
of exact valuation $c$.

Now expand the displacement:
\[
 f(n+T)-f(n)=T\Delta_1f(n)
          +\sum_{k=2}^{T}\binom Tk\Delta_1^kf(n).
\]
For $1\le k\le T$, the binomial denominator identity implies
\[
 v_2\binom Tk\ge v_2(T)-v_2(k)
            \ge v_2(T)-\lfloor\log_2k\rfloor.
\]
Every term of the remaining sum consequently has valuation at least
$v_2(T)+c+1$.  Its sum cannot cancel the leading term, which has exact
valuation $v_2(T)+c$.  This proves both nonzeroness and equality of
valuations.  Only finite sums have been used.
The general theorem is
\lean{KomorebiPaperIX.NewtonDominance.exact_displacement}.
\end{proof}

\begin{theorem}[Exact four-packet displacement law]\label{thm:arith-exact}
Set
\[
 (c_0,c_1,c_2,c_3)=(1,1,4,2),\qquad
 (d_0,d_1,d_2,d_3)=(1,1,5,3).
\]
For $0\le j<4$, $n\ge0$, and every positive integer $T$,
\begin{align*}
 C_j(n+T)-C_j(n)&\ne0,&
 v_2\bigl(C_j(n+T)-C_j(n)\bigr)&=v_2(T)+c_j,\\
 U_j(n+T)-U_j(n)&\ne0,&
 v_2\bigl(U_j(n+T)-U_j(n)\bigr)&=v_2(T)+d_j.
\end{align*}
\end{theorem}

\begin{proof}
The first Newton coefficients, derived from
\eqref{eq:arith-triangular}, are
\[
 a_{0,1}=-2,\quad a_{1,1}=-\frac65,\quad
 a_{2,1}=-\frac{16}{15},\quad a_{3,1}=-\frac{36}{35}.
\]
Their exact valuations are $c_j$.  Theorem~\ref{thm:arith-dyadic-differences}
gives $v_2(a_{j,k})\ge k$ for all orders.  To obtain the stronger
logarithmic envelope required by Lemma~\ref{lem:arith-dominance}, only
a finite initial segment needs improvement.  The recurrence proves the
following exact valuations:
\begin{center}
\begin{tabular}{cll}\toprule
Address $j$ & Further orders $k$ & $v_2(a_{j,k})$ in that order\\\midrule
$0$ & $2$ & $3$\\
$1$ & $2$ & $3$\\
$2$ & $2,3,4,5,6$ & $6,9,13,16,22$\\
$3$ & $2,3,4$ & $8,8,15$\\\bottomrule
\end{tabular}
\end{center}
For all subsequent $k$, the elementary inequality
$c_j+1+\lfloor\log_2k\rfloor\le k$ completes the envelope.
The values in the table are exact rational computations from the defining
recurrence, with their valuations proved in Lean.  They are not a
finite sample used to infer an unproved general pattern.
The uniform envelope is
\lean{KomorebiPaperIX.DyadicFourPrefix.tail_precision}.

Apply Lemma~\ref{lem:arith-dominance} to $C_j$.  Multiplication by the
nonzero factorial $j!$ adds $v_2(j!)$ and preserves nonzeroness.
Hence $d_j=c_j+v_2(j!)$, giving the stated four original gains.
The declarations are
\lean{KomorebiPaperIX.DyadicFourLeastPeriods.normalized_exact_displacement}
and \lean{KomorebiPaperIX.DyadicFourLeastPeriods.exact_displacement}.
\end{proof}

\begin{corollary}[All positive periods and their least element]
\label{cor:arith-least-period}
For $r\in\N$ and $0\le j<4$, a positive integer $T$ is a packet period
of $U_j$ modulo $2^r$ if and only if
\[
 2^{\max(0,r-d_j)}\mid T.
\]
Here a packet period means $U_j(n+T)\equiv U_j(n)\pmod{2^r}$ for every
$n\ge0$.  Consequently the least positive packet periods are
\[
 \begin{array}{c|cccc}
 j&0&1&2&3\\\hline
 \text{least period}&2^{\max(0,r-1)}&2^{\max(0,r-1)}&
                 2^{\max(0,r-5)}&2^{\max(0,r-3)}.
 \end{array}
\]
\end{corollary}

\begin{proof}
The exact displacement law says that the congruence at any row is
 equivalent to $v_2(T)+d_j\ge r$.  This is exactly the displayed integer
divisibility criterion.  Its least positive solution is the indicated
power of two, which also covers $r=0$ and $r\le d_j$.
The all-period and least-element statements are
\lean{KomorebiPaperIX.DyadicFourLeastPeriods.period_iff_dvd} and
\lean{KomorebiPaperIX.DyadicFourLeastPeriods.least_packet_period}.
\end{proof}

For example, $U_2(n)\equiv2\pmod{32}$ and $U_3(n)\equiv6\pmod8$
for every row.  At precision $2^{10}$ the least packet periods are
$512,512,32,128$.  A packet displacement $T$ has original degree
displacement $4T$, as proved by
\lean{KomorebiPaperIX.DyadicFourLeastPeriods.original_degree_exact_displacement}.
This correspondence does not assert minimality among every possible
shift of the full degree-indexed sequence with its off-support zero
coefficients.  Exact valuations of arbitrary higher-order differences
are also a separate question from the first-displacement law.

\subsection{Classical odd-length coordinates and their congruence ladder}

The classical signed Euler numbers have generating function
$\operatorname{sech}t=\sum_{n\ge0}E_{2n}t^{2n}/(2n)!$.
They are the $N=2$, $j=0$ member of
\eqref{eq:arith-euler-egf}.  The original coordinate theorem reconstructs
them from finite alternating sums of odd powers.

\begin{definition}[Odd-length source and inverse coefficients]
\label{def:arith-odd-source}
For a positive odd integer $p$, define
\[
 S_d(p)=\sum_{a=0}^{p-1}(-1)^a(2a+1)^d.
\]
Let $\gamma_m$ be the integral Hurwitz coefficients of $\operatorname{sech}^2$:
\[
 \operatorname{sech}^2t=\sum_{m\ge0}\gamma_m\frac{t^{2m}}{(2m)!},
 \qquad
 \gamma_m=\sum_{b=0}^{2m}\binom{2m}{b}E_bE_{2m-b},
\]
where $E_b=0$ for odd $b$.  Then $\gamma_0=1$, and define
\[
 T^{\mathrm{odd}}_{p,M}(n)=\sum_{m=0}^{\min(M,n)}
       \binom{2n}{2m}\gamma_m p^{2m}S_{2(n-m)}(p).
\]
Here $p$ need only be odd; primality is needed only to interpret the
result as a prime-adic congruence ladder.
\end{definition}

The source and inverse coefficients are
\lean{KomorebiPaperIX.OddPower.source} and
\lean{KomorebiPaperIX.EulerCoefficients.inverseEven}; the truncation is
\lean{KomorebiPaperIX.OddEuler.truncation}.

\begin{theorem}[Classical odd-length reconstruction]
\label{thm:arith-odd-reconstruction}
For every positive odd integer $p$ and every $n\in\N$,
\[
 E_{2n}=\sum_{m=0}^{n}\binom{2n}{2m}\gamma_m p^{2m}S_{2(n-m)}(p).
\]
Equivalently, $E_{2n}=\sum_{q=0}^n V_{nq}S_{2q}(p)$ for the lower
triangular matrix
\[
 V_{nq}=\binom{2n}{2q}\gamma_{n-q}p^{2(n-q)}\quad(q\le n),
 \qquad V_{nq}=0\quad(q>n).
\]
It has $V_{nn}=1$ and $p^{2(n-q)}\mid V_{nq}$ when $q\le n$.
\end{theorem}

\begin{proof}
The literal exponential generating series of the source is the finite sum
\[
 \mathcal S_p(t)=\sum_{a=0}^{p-1}(-1)^a e^{(2a+1)t}.
\]
Multiplication by $\cosh t$ telescopes.  Since $p$ is odd,
\[
 \mathcal S_p(t)=\frac{1+e^{2pt}}{2}\operatorname{sech}t.
\]
Taking its even part gives
\[
 \sum_{n\ge0}S_{2n}(p)\frac{t^{2n}}{(2n)!}
       =\cosh^2(pt)\operatorname{sech}t.
\]
Multiplication by the actual reciprocal $\operatorname{sech}^2(pt)$
yields $\operatorname{sech}t$.  The EGF product rule now gives the
reconstruction coefficient by coefficient.  Replacing $m$ by $n-q$
produces the displayed lower triangular row.  Its diagonal and weights
follow from $\gamma_0=1$ and the explicit factor $p^{2(n-q)}$.
The source-series identity is
\lean{KomorebiPaperIX.OddPower.odd_source_even_egf}; the reconstruction
and row statement are
\lean{KomorebiPaperIX.OddEuler.odd_euler_reconstruction} and
\lean{KomorebiPaperIX.OddEuler.row_action}.
\end{proof}

\begin{corollary}[Odd-prime truncation ladder]\label{cor:arith-odd-ladder}
For every positive odd integer $p$ and all $M,n\in\N$,
\[
 p^{2M+2}\mid E_{2n}-T^{\mathrm{odd}}_{p,M}(n)\quad\text{in }\Z.
\]
If $M\ge n$, the truncation equals $E_{2n}$ exactly.  In particular,
for an odd prime $p$ this is a congruence modulo $p^{2M+2}$ in $\Z_p$.
\end{corollary}

\begin{proof}
Every omitted summand has $m\ge M+1$, hence contains the integer factor
$p^{2m}$ divisible by $p^{2M+2}$.  All its other factors are integral.
The finite sum of these omitted terms proves the assertion, with the
empty-tail case giving exact equality.  These are
\lean{KomorebiPaperIX.OddEuler.odd_euler_supercongruence} and
\lean{KomorebiPaperIX.OddEuler.full_depth_exact}.
\end{proof}

This concrete coordinate theorem is valid before passage to a completion.
The dyadic coordinate source is different.  Its original bound of
$2^{2M+2}$ follows from integral reciprocal coefficients, and the next
subsection sharpens that bound to the optimal uniform exponent $4M+3$.

\subsection{The optimal classical dyadic edge tower}

Return now to the classical signed Euler numbers
\[
 \operatorname{sech}t=\frac1{\cosh t}
       =\sum_{n\ge0}E_{2n}\frac{t^{2n}}{(2n)!}.
\]
Thus $E_{2n}=E_{2n}^{(2,0)}$ in the convention
\eqref{eq:arith-euler-egf}.  In particular $E_0=1$, $E_2=-1$, and $E_4=5$.
The following construction concerns these $N=2$ coefficients, not $C_{4,0}$.

\begin{definition}[Literal dyadic truncation]\label{def:arith-truncation}
Define $\beta_m$ by the formal unit-series identity
\[
 \sum_{m\ge0}\beta_m z^{2m}=\left(\frac{z}{\sinh z}\right)^2,
 \qquad \beta_0=1,
\]
where the quotient means the inverse of the unit series
$(\sinh z/z)^2$.  Define
\[
 A(n)=\frac{3^{2n+2}-1}{4(2n+1)(2n+2)},\qquad
 S_m(n)=\beta_m2^{2m}\frac{(2n)!}{(2n-2m)!}A(n-m)
       \quad(0\le m\le n),
\]
and
\begin{equation}\label{eq:arith-tm}
 T_M(n)=\sum_{m=0}^{\min(M,n)}S_m(n)\qquad(M,n\in\N).
\end{equation}
\end{definition}

These are the actual source definitions
\lean{KomorebiPaperIX.Dyadic.beta},
\lean{KomorebiPaperIX.Dyadic.sourceEven}, and
\lean{KomorebiPaperIX.DyadicReconstruction.truncation}.
The formal summand uses a falling factorial, and
\lean{KomorebiPaperIX.DyadicReconstruction.falling_factorial_fraction}
proves its equality with the quotient displayed above.

\begin{lemma}[Exact reconstruction and source integrality]
\label{lem:arith-reconstruction}
For every $n\in\N$,
\[
 E_{2n}=\sum_{m=0}^{n}S_m(n).
\]
Every $\beta_m$ is $2$-adically integral and every $A(n)$ is a nonzero
$2$-adic unit.  If $M\ge n$, then $T_M(n)=E_{2n}$.
\end{lemma}

\begin{proof}
The ordinary kernel coefficients of $(\sinh z/z)^2$, in $z^2$, are
\[
 \frac{2^{2m+1}}{(2m+2)!}.
\]
Legendre's strict factorial bound
$v_2((2m+2)!)<2m+2$ makes them integral.  Their constant coefficient is
one.  The coefficient recurrence of the genuine reciprocal therefore
proves integrality of every $\beta_m$.
For the other factor, lifting the exponent at the even exponent $2n+2$
gives
\[
 v_2(3^{2n+2}-1)=2+v_2(2n+2).
\]
Its denominator has the same valuation, since $2n+1$ is odd.  Thus
$A(n)$ is a nonzero unit.

Its exponential generating function is
\[
 \mathcal A(t)=\sum_{n\ge0}A(n)\frac{t^{2n}}{(2n)!}
             =\frac{\cosh(3t)-\cosh t}{4t^2}.
\]
All quotients by powers of $t$ here remove proved initial zero
coefficients.  The hyperbolic identity
\[
 \mathcal A(t)\cosh t=\left(\frac{\sinh(2t)}{2t}\right)^2
\]
then gives the exact formal-series product
\[
 \left(\sum_{m\ge0}\beta_m(2t)^{2m}\right)\mathcal A(t)
       =\operatorname{sech}t.
\]
Taking the factorial-scaled coefficient of $t^{2n}$ produces the stated
finite sum.  Full depth leaves no omitted terms.
The relevant declarations are
\lean{KomorebiPaperIX.Dyadic.beta_integral},
\lean{KomorebiPaperIX.Dyadic.source_valuation},
\lean{KomorebiPaperIX.DyadicSeries.dyadic_reconstruction_series}, and
\lean{KomorebiPaperIX.DyadicReconstruction.dyadic_reconstruction}.
\end{proof}

\begin{lemma}[Exact weighted inverse coefficient]\label{lem:arith-kappa}
Let $\kappa_m=\beta_m2^{2m}(2m)!$.  For every $m\ge1$,
\[
 \kappa_m=-2^{4m}(2m-1)B_{2m}\ne0,\qquad
 v_2(\kappa_m)=4m-1,
\]
where $B_{2m}$ is the usual even Bernoulli number.
\end{lemma}

\begin{proof}
Let $Q(z)=z\coth z$, the regular formal series with constant coefficient
one.  The elementary differential identity
\[
 \left(\frac z{\sinh z}\right)^2=Q(z)-zQ'(z)
\]
and the Bernoulli expansion
\[
 Q(z)=1+\sum_{m\ge1}\frac{2^{2m}B_{2m}}{(2m)!}z^{2m}
\]
give
$\beta_m=-2^{2m}(2m-1)B_{2m}/(2m)!$.
Multiplication by $2^{2m}(2m)!$ proves the identity for $\kappa_m$.
The pinned Bernoulli valuation theorem gives $v_2(B_{2m})=-1$;
since $2m-1$ is odd, the valuation is exactly $4m-1$.
This also proves nonzeroness.  The source proves the formal hyperbolic
identity after clearing removable zeros and then invokes the library
Bernoulli theorem, rather than assuming the value of $\kappa_m$.
The declarations are
\lean{KomorebiPaperIX.DyadicSharpTower.kappa_bernoulli} and
\lean{KomorebiPaperIX.DyadicSharpTower.kappa_valuation}.
\end{proof}

\begin{theorem}[Optimal uniform dyadic tower]\label{thm:arith-sharp-tower}
For all $M,n\in\N$,
\begin{equation}\label{eq:arith-sharp-tail}
 E_{2n}-T_M(n)\in2^{4M+3}\Z_2.
\end{equation}
At the first omitted row $n=M+1$, the error is nonzero and
\[
 v_2\bigl(E_{2M+2}-T_M(M+1)\bigr)=4M+3.
\]
Thus the exponent $4M+3$ cannot be increased uniformly in the Euler row.
\end{theorem}

\begin{proof}
Rewrite the summand exactly as
\[
 S_m(n)=\kappa_m\binom{2n}{2m}A(n-m).
\]
For $m>M$, its first factor is divisible by $2^{4m-1}$, the binomial
coefficient is integral, and $A(n-m)$ is a unit.  Hence every omitted
summand is divisible by $2^{4M+3}$.  Summing the finite tail proves
\eqref{eq:arith-sharp-tail}.  If $M\ge n$, the tail is empty and the
statement includes the resulting zero error.

At $n=M+1$ there is exactly one omitted summand.  Here
$\binom{2n}{2n}=1$ and $A(0)=1$, so the error is $\kappa_{M+1}$.
Lemma~\ref{lem:arith-kappa} gives its exact valuation $4M+3$ and
nonzeroness.  The statements are
\lean{KomorebiPaperIX.DyadicSharpTower.dyadic_supercongruence_sharp},
\lean{KomorebiPaperIX.DyadicSharpTower.first_omitted_error}, and
\lean{KomorebiPaperIX.DyadicSharpTower.first_omitted_error_valuation}.
\end{proof}

\begin{proposition}[Row-dependent summand weights]\label{prop:arith-row-weight}
For $1\le m\le n$, the summand is nonzero and
\[
 v_2(S_m(n))=4m-1+v_2\binom{2n}{2m}.
\]
More generally, if $L\in\N$ satisfies
\[
 L\le4m-1+v_2\binom{2n}{2m}\qquad\text{for every }M<m\le n,
\]
then $E_{2n}-T_M(n)\in2^L\Z_2$.
\end{proposition}

\begin{proof}
Use the same factorization of $S_m(n)$, the nonzero binomial coefficient
for $m\le n$, the exact valuation of $\kappa_m$, and the unit property
of $A(n-m)$.  The tail conclusion follows by addition in $2^L\Z_2$.
These are \lean{KomorebiPaperIX.DyadicSharpTower.summand_valuation} and
\lean{KomorebiPaperIX.DyadicSharpTower.dyadic_supercongruence_row_weight}.
\end{proof}

The minimum of the summand weights is a lower bound for the valuation of
the tail.  Cancellation can make the tail more divisible.  Thus this
proposition gives exact individual weights, while an exact formula for
every row's entire tail remains a distinct problem.  It does not weaken
the uniform optimality already established at $n=M+1$.

\subsection{Isometric continuation of the four dyadic packets}

The exact displacement law has a stronger consequence than continuity of
the original packets.  It identifies their completed images exactly.

\begin{theorem}[Isometric packet bijections]\label{thm:arith-isometry}
For $0\le j<4$, define on natural numbers
\[
 f_j(n)=\frac{U_j(n)-j!}{2^{d_j}},\qquad (d_0,d_1,d_2,d_3)=(1,1,5,3).
\]
There is a unique continuous map $\Phi_j:\Z_2\to\Z_2$ whose value at
every $n\in\N$ equals the embedded rational $f_j(n)$.
It is an isometric bijection and $\Phi_j(0)=0$:
\[
 |\Phi_j(x)-\Phi_j(y)|_2=|x-y|_2\qquad(x,y\in\Z_2).
\]
In particular, every element of $\Z_2$ has exactly one preimage.
\end{theorem}

\begin{proof}
For $n>0$, the displacement law at row zero gives
$v_2(U_j(n)-j!)=v_2(n)+d_j$, so $f_j(n)\in\Z_2$.
The value at zero is zero.  Applying the same law to the smaller of two
unequal natural numbers shows
\[
 |f_j(n)-f_j(m)|_2=|n-m|_2.
\]
The natural numbers are dense in $\Z_2$, and the target is complete.
The map therefore extends continuously.  The distance identity extends
on the dense product of natural-number points, making the extension
an isometry.  Density also proves uniqueness among all continuous maps
with these prescribed samples.

Surjectivity uses compactness, and is proved separately rather than
inferred merely from injectivity.  For an isometry $g$ of a compact
metric space into itself, the sequence of iterates $g^n(x)$ has a
convergent subsequence.  Two sufficiently close iterates, with exponents
$a<b$, yield $d(x,g^{b-a}(x))$ equally small after removing $a$ isometric
iterations.  Since $b-a>0$, the latter point is in the image of $g$.
Thus the image is dense.  It is also compact and hence closed, so it is
the whole space.  Apply this to $\Phi_j$ on $\Z_2$.

The general ingredients are
\lean{KomorebiPaperIX.DensePadicIsometry.existsUnique_extensionEquiv} and
\lean{KomorebiPaperIX.CompactSelfIsometry.surjective}.
The literal packet specialization is
\lean{KomorebiPaperIX.DyadicFourIsometry.existsUnique_isometry};
uniqueness among continuous maps is
\lean{KomorebiPaperIX.DyadicFourIsometry.unique_continuous}.
\end{proof}

\begin{corollary}[Exact image of the original coefficients]
\label{cor:arith-affine-image}
The continuous extension of the original sequence is
\[
 \widetilde U_j(x)=j!+2^{d_j}\Phi_j(x).
\]
It is injective, its image is exactly $j!+2^{d_j}\Z_2$, and
\[
 |\widetilde U_j(x)-\widetilde U_j(y)|_2
     =2^{-d_j}|x-y|_2.
\]
Every point of that affine image has exactly one preimage in $\Z_2$.
\end{corollary}

\begin{proof}
The formula agrees with $U_j$ on every natural number.  Continuity and
the distance equality follow from the corresponding properties of
$\Phi_j$.  Its surjectivity gives precisely the displayed affine image,
and multiplication by the nonzero scalar $2^{d_j}$ preserves injectivity.
The declarations are
\lean{KomorebiPaperIX.DyadicFourIsometry.originalExtension_nat},
\lean{KomorebiPaperIX.DyadicFourIsometry.originalExtension_dist}, and
\lean{KomorebiPaperIX.DyadicFourIsometry.originalExtension_range}.
\end{proof}

This closes the passage from the rational four-packet sequences to
bijections on their exact dyadic images.  The construction uses density,
completeness, and the exact displacement law.  It does not assert
analyticity of these particular packet extensions merely because they
are isometries.  The polynomial analytic solution maps established later
have their own hypotheses and proof.

\section{Finite obstructions, compatible lifts, and deformation}
\label{sec:lifting}

The arithmetic results above concern particular families. This section isolates
the finite correction problem and follows it through polynomial lifting,
continuous parameter charts, and quantitative deformation. The statements keep
track of three different objects: a residual obstruction, the set of correcting
digits, and an actual root in a completion. A vanishing obstruction at one
precision supplies a next lift. Additional hypotheses are needed to continue
indefinitely.

Unless a different coefficient ring is specified, a polynomial system is a
finite list of integer polynomials. All divisibility conditions on vectors are
coordinatewise. For a positive integer $q$, write $S_q=\Z/q\Z$ and use an
overline for reduction modulo $q$. For a prime $p$, write
$\F_p=S_p$, normalize $|p|_p=p^{-1}$, and put
\[
 B_L(x)=x+p^L\Zp^d.
\]
An integer center is included in this notation by its natural image in
$\Zp^d$. All dimensions below are finite and may be zero. A statement over
$S_q$ allows composite $q$ and $q=1$ when explicitly stated. A cardinality
formula in terms of vector-space rank uses prime $p$.

\subsection{The actual polynomial correction equation}

For $F=(F_1,\ldots,F_e)\in\Z[X_1,\ldots,X_d]^e$, define its evaluated
Jacobian by the formal partial derivatives:
\[
 DF(x)_{ij}=(\partial_jF_i)(x),\qquad
 J_x=DF(x)\bmod q:S_q^d\longrightarrow S_q^e.
\]
Thus the linear map used in the obstruction is computed from the input
polynomials. The following algebraic Taylor statement supplies the connection
to their actual evaluation.

\begin{lemma}[Polynomial Taylor remainder]\label{lem:polynomial-taylor}
Let $R$ be a commutative ring, $P\in R[X_1,\ldots,X_d]$, $a\in R^d$, and
$c\in R$. There is a polynomial $Q\in R[U_1,\ldots,U_d]$, with $Q(0)=0$,
such that for every $u\in R^d$,
\begin{equation}\label{eq:polynomial-taylor}
 P(a+cu)=P(a)+c\sum_{j=1}^d(\partial_jP)(a)u_j+c^2Q(u).
\end{equation}
In particular, for integer $a,u,c$, the difference between the first two
terms and $P(a+cu)$ is divisible by $c^2$.
\end{lemma}

\begin{proof}
Apply polynomial induction using constants, addition, and multiplication by a
variable. Constants have zero remainder, and remainders add. For the
multiplication step, suppose the assertion is known for $P$ and write
$D_P(U)=\sum_j(\partial_jP)(a)U_j$. Multiplication of its expansion by
$a_i+cU_i$ gives constant term $P(a)a_i$, linear term
$c(a_iD_P(U)+P(a)U_i)$, and remainder
\[
 c^2\bigl(Q(U)(a_i+cU_i)+D_P(U)U_i\bigr).
\]
The product rule identifies the linear term with the actual derivative of
$PX_i$. The displayed remainder vanishes at $U=0$. This constructs one
polynomial remainder before evaluating $u$, including over rings where
division by $c$ would be invalid.
\end{proof}

The general polynomial identity is formalized by
\lean{Komorebi.PolynomialTaylorRing.exists_eval_scaled_remainder}.
The earlier integer divisibility interface is
\lean{Komorebi.PolynomialTaylor.taylor_dvd}.

\begin{theorem}[Finite polynomial obstruction]\label{thm:finite-polynomial-obstruction}
Let $q\ge1$, $L\ge1$, $x\in\Z^d$, and suppose
$F(x)\in q^L\Z^e$. Define the exact integer residual $r_x$ by
$F(x)=q^Lr_x$, and put
\[
 \mathcal O_x=S_q^e/\im J_x,\qquad o_x=[\overline{r_x}].
\]
Then $\mathcal O_x$ is finite and
\begin{equation}\label{eq:finite-obstruction}
 o_x=0\quad\Longleftrightarrow\quad
 \exists y\in\Z^d:\quad
 y\equiv x\pmod{q^L},\qquad F(y)\equiv0\pmod{q^{L+1}}.
\end{equation}
No surjectivity or nonsingularity assumption is required.
\end{theorem}

\begin{proof}
Every point in the lower residue class has the form $y=x+q^Lt$ with
$t\in\Z^d$. By Lemma~\ref{lem:polynomial-taylor} and $2L\ge L+1$,
\[
 F(x+q^Lt)\equiv q^L\bigl(r_x+DF(x)t\bigr)\pmod{q^{L+1}}.
\]
The integer $q^L$ is nonzero. Cancelling it in the corresponding divisibility
condition proves
\begin{equation}\label{eq:digit-correction}
 F(x+q^Lt)\equiv0\pmod{q^{L+1}}
 \quad\Longleftrightarrow\quad
 \overline{r_x}+J_x\overline t=0.
\end{equation}
Integer vectors reduce onto $S_q^d$. The equation on the right is solvable
exactly when $\overline{r_x}$ belongs to $\im J_x$, since that image is
closed under negation. This is precisely $o_x=0$. Finiteness follows because
$\mathcal O_x$ is a quotient of the finite module $S_q^e$.
\end{proof}

The source theorem
\lean{Komorebi.PolynomialObstruction.finite_polynomial_obstruction}
presents the image using the integer-domain map
$\Z^d\xrightarrow{DF(x)}\Z^e\to S_q^e$. Reduction onto $S_q^d$ identifies
that image with $\im J_x$. The class-preserving comparison is explicit in
\lean{Komorebi.PolynomialLiftCounts.obstructionEquiv}; the change of
presentation does not change the distinguished residual class.

\begin{proposition}[Affine depth zero]\label{prop:affine-obstruction}
Let $F(x)=Ax-b$ for an integer linear map $A:\Z^d\to\Z^e$. The criterion
of Theorem~\ref{thm:finite-polynomial-obstruction} holds for every $L\ge0$,
with $J=A\bmod q$.
\end{proposition}

\begin{proof}
The affine identity
$F(x+q^Lt)=F(x)+q^LAt$ has no quadratic remainder. The proof of
\eqref{eq:digit-correction} therefore needs no inequality involving $2L$.
The residual and the finite quotient are obtained as above.
\end{proof}

The exact finite statement is
\lean{Komorebi.NativeAffine.finite_affine_obstruction}.
Positive depth is essential in the nonlinear theorem: for $F(X)=X^2-1$,
$q=2$, $L=0$, and $x=0$, the Jacobian is zero and the residual class is
nonzero, yet $y=1$ is a root modulo two in the same residue modulo one.

\begin{proposition}[Independence of the integer representative]
\label{prop:representative-independence}
Under the hypotheses of Theorem~\ref{thm:finite-polynomial-obstruction},
let $x'=x+q^Ls$ and write $F(x')=q^Lr_{x'}$. Then
\[
 J_{x'}=J_x,\qquad
 \overline{r_{x'}}=\overline{r_x}+J_x\overline s.
\]
Consequently the two residual classes coincide in the common quotient
$S_q^e/\im J_x$.
\end{proposition}

\begin{proof}
Polynomial evaluation respects coordinatewise congruence. Apply this first
to every formal partial derivative: $x'\equiv x\pmod q$ gives
$J_{x'}=J_x$. Next apply \eqref{eq:polynomial-taylor} to $F$ at scale $q^L$.
Subtract the two exact residual identities and cancel $q^L$ to obtain the
displayed relation modulo $q$. Its additional term belongs to the common
Jacobian image, so it vanishes in the quotient.
\end{proof}

The comparison is
\lean{Komorebi.PolynomialRepresentative.representative_independence}.
This proves independence of the integer representative within a fixed
polynomial presentation. It makes no identification of the raw cokernels
of different equation lists. For example, adjoining a zero equation leaves
the roots unchanged while adding an $S_q$ summand to the raw cokernel.

\subsection{One correction shared by two systems}

The next obstruction concerns compatibility between individually attainable
corrections. Its algebraic form works over an arbitrary commutative ring.
Let $U,V,W$ be $R$-modules and let $J:U\to V$, $K:U\to W$ be linear.
Define
\begin{equation}\label{eq:common-quotient}
 C(J,K)=U/(\ker J+\ker K).
\end{equation}
An attainable value $Ju$ has label $[u]\in C(J,K)$; this is independent of
the choice of its preimage because $\ker J$ is killed in the quotient.
Similarly, $Kv$ has label $[v]$. Denote these maps by
$\lambda_J:\im J\to C(J,K)$ and $\lambda_K:\im K\to C(J,K)$.

\begin{theorem}[Common correction and fiber product]\label{thm:common-correction}
For every $u,v\in U$,
\begin{equation}\label{eq:common-correction}
 [u]=[v]\text{ in }C(J,K)
 \quad\Longleftrightarrow\quad
 \exists w\in U:\ Jw=Ju,\quad Kw=Kv.
\end{equation}
Both label maps are surjective, and the image of the joint map is exactly
the fiber product:
\begin{equation}\label{eq:joint-fiber-product}
 \im\bigl(U\to\im J\times\im K,\ u\mapsto(Ju,Ku)\bigr)
 =\im J\mathbin{\times}_{C(J,K)}\im K.
\end{equation}
If either $V$ or $W$ is finite, then $C(J,K)$ is finite.
\end{theorem}

\begin{proof}
If $[u]=[v]$, write $u-v=s+t$ with $s\in\ker J$ and $t\in\ker K$.
Then $w=u-s=v+t$ satisfies both required equalities. Conversely, a common
$w$ gives
\[
 u-v=(u-w)+(w-v)\in\ker J+\ker K.
\]
Every class $[u]$ is the label of $Ju$ and of $Ku$, proving surjectivity.
For attainable $a=Ju$ and $b=Kv$, the first equivalence says exactly that
equal labels are equivalent to a simultaneous preimage. This proves
\eqref{eq:joint-fiber-product}. Finally, a finite target has finite image,
and the associated surjective label map makes $C(J,K)$ finite.
\end{proof}

These statements are
\lean{Komorebi.SharedCorrection.quotient_mk_eq_iff_common} and
\lean{Komorebi.SharedCorrection.joint_range_eq_fiber}, together with the
finite-target corollaries in the same module. The quotient also has the
expected universal property: every linear map from $U$ annihilating both
kernels factors uniquely through $C(J,K)$. Equivalently, compatible linear
observations on $\im J$ and $\im K$ factor uniquely through their label
maps. This is formalized by
\lean{Komorebi.SharedCorrection.common_paired_universal}.

\begin{corollary}[Simultaneous polynomial lifting]\label{cor:shared-polynomial}
Let $F\in\Z[X_1,\ldots,X_d]^e$ and
$G\in\Z[X_1,\ldots,X_d]^f$ have a common lower root $x$ modulo $q^L$,
where $q\ge1$ and $L\ge1$. Use the integer-domain reduced Jacobians
$J:\Z^d\to S_q^e$ and $K:\Z^d\to S_q^f$. Suppose $u,v\in\Z^d$ are
separately valid corrections for $F$ and $G$, respectively. Then
\[
 [u]=[v]\text{ in }\Z^d/(\ker J+\ker K)
\]
if and only if a single $y\equiv x\pmod{q^L}$ solves both systems modulo
$q^{L+1}$. Each separate label is independent of the chosen valid
correction. The common quotient is finite.
\end{corollary}

\begin{proof}
For each system, \eqref{eq:digit-correction} translates a valid next root
into a linear correction equation. Once $u$ corrects $F$, another digit
$w$ corrects $F$ precisely when $Jw=Ju$. The analogous condition for $G$
is $Kw=Kv$. Theorem~\ref{thm:common-correction} now gives the criterion.
Two valid corrections of one block differ by a vector in its Jacobian
kernel, so they have equal labels. Finiteness follows from either finite
residue target. An unconditional formulation quantifies over separately
valid $u,v$ with equal labels, thereby including their existence in the
criterion rather than assuming it.
\end{proof}

See \lean{Komorebi.PolynomialShared.joint_lift_iff_class_eq} and
\lean{Komorebi.PolynomialShared.finite_polynomial_pair_compatibility}.
Separate liftability alone is insufficient. At $x=0$ modulo two,
$F(X)=X$ and $G(X)=X-2$ both lift to modulus four, but their correcting
digits are respectively $0$ and $1$ modulo two. No single residue solves
both equations modulo four.

\subsection{Correction torsors and exact residue counts}

It is useful to distinguish integer representatives from next-level residue
classes. For $q\ge1$ and $L\ge0$, the map
\[
 S_q^d\longrightarrow
 \{y\in(\Z/q^{L+1}\Z)^d:y\equiv x\pmod{q^L}\},
 \qquad t\longmapsto x+q^L\widetilde t
\]
is a bijection, where $\widetilde t$ is any integer representative.
Independence of the representative and injectivity both follow by cancelling
$q^L$ in integer divisibility; surjectivity follows by expressing the lower
congruence as a multiple of $q^L$. This is the equivalence
\lean{Komorebi.ResidueLiftDigits.digitEquivResidue}.

At positive polynomial depth, restrict it to the correction fiber
\[
 \mathcal T_x=\{t\in S_q^d:\overline{r_x}+J_xt=0\}
\]
and the actual lift set
\[
 \mathcal L_{x,L}=
 \{y\in(\Z/q^{L+1}\Z)^d:y\equiv x\pmod{q^L},\ F(y)=0\}.
\]
The polynomial predicate in the second display is evaluation in the residue
ring. Evaluation on canonical integer representatives gives the same
predicate, because integer polynomial evaluation respects congruence.

\begin{theorem}[Torsor and lift cardinalities]\label{thm:lift-torsor-count}
Under the hypotheses of Theorem~\ref{thm:finite-polynomial-obstruction},
there is an equivalence $\mathcal T_x\simeq\mathcal L_{x,L}$.
Every nonempty fiber is a torsor for $\ker J_x$ under addition of
correction digits. If $q=p$ is prime and
$\rho_x=\rank_{\F_p}J_x$, then
\begin{equation}\label{eq:lift-count}
 \#\mathcal L_{x,L}=
 \begin{cases}p^{d-\rho_x},&o_x=0,\\0,&o_x\ne0,\end{cases}
 \qquad
 \#\mathcal O_x=p^{e-\rho_x}.
\end{equation}
\end{theorem}

\begin{proof}
Equation~\eqref{eq:digit-correction} identifies exactly which points of the
residue fiber satisfy the equations, proving the restricted equivalence.
If $u,t\in\mathcal T_x$, subtraction gives $J_x(t-u)=0$. Conversely,
adding any element of $\ker J_x$ preserves the correction equation. The
unique kernel vector carrying $u$ to $t$ is $t-u$, so the action is free
and transitive. Choosing one valid digit identifies the fiber with the
kernel, without making that choice part of its translation action.

For prime $p$, rank-nullity gives
$\dim\ker J_x=d-\rho_x$ and
$\dim(\F_p^e/\im J_x)=e-\rho_x$. The cardinality of an
$n$-dimensional $\F_p$-space is $p^n$. Combine these identities with the
obstruction criterion and the residue equivalence.
\end{proof}

The main count is
\lean{Komorebi.PolynomialLiftCounts.finite_polynomial_lift_count}.
The explicit torsor and chosen-point equivalence are developed in
\lean{Komorebi.FiniteFiberCounts.translate_exists_unique} and
\lean{Komorebi.PolynomialLiftCounts.liftEquivKernel}.
For composite $q$, the torsor and kernel cardinality remain valid, while the
vector-space rank formula is unavailable. Multiplication by two on
$\Z/4\Z$ already has a kernel of size two.

\subsection{Smooth compatible towers and exact roots}

Surjectivity removes every residual obstruction and persists throughout a
positive-depth residue ball. This supplies the additional hypothesis needed
to iterate the finite theorem.

\begin{theorem}[Compatible polynomial tower]\label{thm:smooth-tower}
Let $q\ge1$, $L\ge1$, $F\in\Z[X_1,\ldots,X_d]^e$, and
$x\in\Z^d$ satisfy $F(x)\in q^L\Z^e$. Suppose $J_x:S_q^d\to S_q^e$
is surjective. There is a single sequence $z_n\in\Z^d$ such that
\begin{align*}
 z_0&=x,& F(z_n)&\in q^{L+n}\Z^e,\\
 z_{n+1}-z_n&\in q^{L+n}\Z^d,& J_{z_n}&=J_x.
\end{align*}
For $a\le b$, it also satisfies $z_b-z_a\in q^{L+a}\Z^d$.
\end{theorem}

\begin{proof}
At level $N=L+n$, divide the residual by $q^N$. Surjectivity supplies
$t\in S_q^d$ with $J_{z_n}t=-\overline{F(z_n)/q^N}$. Choose an integer
representative and apply \eqref{eq:digit-correction} to obtain $z_{n+1}$.
The new point is congruent to $z_n$ modulo $q$. Evaluation of each formal
partial derivative therefore gives $J_{z_{n+1}}=J_{z_n}$.

Dependent recursion carries the current root and this fixed Jacobian as
its state. Thus surjectivity at all later stages is a consequence of the
single initial assumption. The cross-level congruence follows by induction,
adding consecutive differences whose divisibilities all imply divisibility
by $q^{L+a}$.
\end{proof}

The precise strong statement is
\lean{Komorebi.PolynomialIntegerTower.exists_compatible_tower_strong}.
The proof constructs a compatible path using classical choice. It does not
claim an optimized procedure for selecting its digits.

\begin{theorem}[Completion and smooth branching]\label{thm:smooth-padic-root}
In Theorem~\ref{thm:smooth-tower}, let $q=p$ be prime. The tower can be
chosen with a limit $a\in\Zp^d$ such that
\[
 F(a)=0,\qquad z_n\longrightarrow a,\qquad
 a-z_n\in p^{L+n}\Zp^d\quad(n\ge0).
\]
At each stage, the number of distinct next-level residue roots over $z_n$
is exactly $p^{d-e}$, and the finite obstruction module has one element.
If $d=e$, there is exactly one root of $F$ in $B_L(x)$.
\end{theorem}

\begin{proof}
Cross-level compatibility bounds the supremum distance between later terms
by $p^{-(L+n)}$. Completeness gives a limit, and closed norm bounds retain
the congruence to every $z_n$. The coefficients of the original integer
polynomials are explicitly included in $\Zp$. Their evaluations satisfy
both $F_i(z_n)\to F_i(a)$, by continuity, and $F_i(z_n)\to0$, by their
growing divisibility. Uniqueness of limits gives exact vanishing.

Surjectivity gives $\rank J_x=e\le d$. Theorem~\ref{thm:lift-torsor-count}
then gives $p^{d-e}$ next residue roots and obstruction cardinality one at
every stage.

For square uniqueness, let $a,b\in B_L(x)$ be arbitrary exact roots.
Their canonical integer truncations are genuine roots in the respective
finite residue rings, since evaluation commutes with reduction. They agree
at depth $L$. If they agree at depth $N\ge L$, both next residues belong
to the same lift set, which has cardinality $p^{d-d}=1$. Induction gives
agreement at every depth. The residue maps separate $p$-adic integers, so
$a=b$. Existence was already supplied by the tower.
\end{proof}

The endpoint construction is
\lean{Komorebi.PolynomialPadicLift.exists_padic_root_with_tower}, the
branching statement is
\lean{Komorebi.PolynomialSmoothCounts.exists_regular_tower}, and uniqueness
for all roots in the ball is
\lean{Komorebi.PolynomialPadicUnique.exists_unique_padic_root}.

\begin{remark}\label{rem:singular-tower-boundary}
The one-step singular criterion and the smooth infinite theorem have
different hypotheses. The polynomial $X^2-5$ has odd roots modulo two and
four but has no root modulo eight. Thus a previous zero obstruction does
not force the next one to vanish. Also, rectangular surjectivity alone
does not imply uniqueness: $F(X,Y)=X$ has the distinct dyadic roots
$(0,0)$ and $(0,2)$ in the same residue ball modulo two. These are explicit
scope controls for the formalized tower theorem.
\end{remark}

\subsection{Coefficient stability and continuous root charts}

\begin{proposition}[Smooth coefficient stability]\label{prop:smooth-coefficient-stability}
Let $p$ be prime, $M\ge1$, $F,H\in\Z[X_1,\ldots,X_d]^e$, and
$c\in p^M\Z$. Suppose $DF(x)\bmod p$ is surjective for an integer
$x\in\Z^d$. If $a\in\Zp^d$ satisfies $F(a)=0$ and
$a\equiv x\pmod p$, then there is a root $b$ of $G=F+cH$ with
$b-a\in p^M\Zp^d$.

If $d=e$ and $M\ge L\ge1$, every pair of roots $a,b\in B_L(x)$ of
$F$ and $G$, respectively, satisfies $b-a\in p^M\Zp^d$. The perturbed
root is unique in that ball whenever it exists.
\end{proposition}

\begin{proof}
Since $p\mid c$, differentiation of the actual polynomial list gives
$DG(x)\equiv DF(x)\pmod p$. Truncate $a$ to an integer vector $y$ at
depth $M$. Then $F(y)\equiv0\pmod{p^M}$, hence also
$G(y)\equiv0\pmod{p^M}$. Congruence of $y$ and $x$ modulo $p$ preserves
the reduced Jacobian. The smooth root theorem gives
$b\equiv y\equiv a\pmod{p^M}$.

In the square case this constructed root lies in $B_L(x)$ because $M\ge L$.
Square uniqueness for $G$ identifies it with every competing root in that
ball. The construction therefore gives the precision assertion for an
arbitrary root there.
\end{proof}

See \lean{Komorebi.PolynomialRootStability.exists_nearby_perturbed_root}
and \lean{Komorebi.PolynomialRootStability.perturbed_roots_congruent}.
The scalar $c$ may be negative or zero. The later weighted theorems supply
an arbitrary integral $p$-adic coefficient interface and an explicit
first-order displacement.

\begin{theorem}[Continuous charts with a selected Jacobian minor]
\label{thm:continuous-parameter-chart}
Let $p$ be prime, $L\ge1$,
$F\in\Z[X_1,\ldots,X_{e+k}]^e$, and $x\in\Z^{e+k}$ satisfy
$F(x)\in p^L\Z^e$. Suppose the full reduced Jacobian is surjective.
There is a permutation $\sigma$ of the coordinates such that, writing
$\sigma x=(y_0,t_0)$ with $e$ solved coordinates and $k$ parameter
coordinates, the parameter ball
$T=t_0+p^L\Zp^k$ is homeomorphic to the entire root locus
\[
 \mathcal R=\{z\in B_L(x):F(z)=0\}.
\]
The inverse homeomorphism is the actual free-coordinate projection
$z\mapsto\operatorname{free}(\sigma z)$. If $h:T\to\mathcal R$ is the
forward map, then for every $N\ge0$,
\begin{equation}\label{eq:chart-congruence}
 h(t)\equiv h(u)\pmod{p^N}
 \quad\Longleftrightarrow\quad t\equiv u\pmod{p^N}.
\end{equation}
Equivalently, $\sigma h(t)=(\phi(t),t)$ for a continuous
$\phi:T\to y_0+p^L\Zp^e$ whose graph is the entire reordered root locus.
\end{theorem}

\begin{proof}
First consider a square integer polynomial map $P$ whose reduced Jacobian
at $x$ is surjective. Its Jacobian is invertible over $\F_p$. A digit
comparison using Taylor expansion shows that equal outputs modulo
$p^{N+1}$ force equal inputs modulo $p^{N+1}$ whenever the inputs already
agree modulo $p^N$ and belong to $B_1(x)$. Induction and integer truncation
give, for arbitrary $a,b\in B_1(x)$,
\[
 P(a)\equiv P(b)\pmod{p^N}
 \quad\Longleftrightarrow\quad a\equiv b\pmod{p^N}.
\]

More generally, for a rectangular map with surjective reduced Jacobian,
\begin{equation}\label{eq:polynomial-ball-image}
 P(B_L(x))=P(x)+p^L\Zp^e.
\end{equation}
One inclusion follows from polynomial congruence. For integer targets in
the right-hand ball, subtract the target from the constant terms and use
the smooth root theorem. For arbitrary $p$-adic targets, approximate by
integer truncations. The input ball is compact, so its continuous image is
closed. It contains all these approximants and hence their limit. This
argument does not assume that independently chosen preimages form a
compatible sequence.

For square $P$, the resulting bijection between the two compact residue
balls is continuous, and its target is Hausdorff. It is therefore a
homeomorphism, with the same congruence preservation for its inverse.

Now suppose a chosen set of $e$ columns of $DF(x)\bmod p$ is independent,
and put those coordinates first. Apply the square result to the actual
augmented polynomial map
\[
 \Psi(Y,T)=(F(Y,T),T).
\]
Its reduced Jacobian has upper blocks $(D_YF,D_TF)$ and lower blocks
$(0,I_k)$, so it is invertible. For $t\in T$, the target $(0,t)$ belongs
to $\Psi(x)+p^L\Zp^{e+k}$ because $F(x)\equiv0\pmod{p^L}$.
The inverse of $\Psi$ at that target gives the unique root with parameters
$t$. The inverse congruence property gives one direction of
\eqref{eq:chart-congruence}; projection onto the actual parameter
coordinates gives the other.

Finally, surjectivity of the original rectangular Jacobian guarantees a
basis among its columns. Extend the corresponding injection of coordinate
indices to a permutation $\sigma$. Renaming the variables by
$\sigma^{-1}$ transports polynomial evaluation and every formal partial
derivative. The fixed-coordinate construction therefore applies to the
original system after undoing the permutation.
\end{proof}

The main theorem is
\lean{Komorebi.PolynomialGeneralChart.exists_parameter_homeomorph}.
The actual square inverse and arbitrary-target surjectivity are
\lean{Komorebi.PolynomialLocalInverse.localHomeomorph} and
\lean{Komorebi.PolynomialLocalSurjective.exists_padic_preimage}.
The center $x$ is only an approximate root: the chart need not satisfy
$\phi(t_0)=y_0$ unless $F(x)=0$ exactly. The assertion selects suitable
coordinates; it does not permit an arbitrary prescribed solved block.

\subsection{Prescribed parameters and their obstruction}

Suppose now that the split into solved variables and parameter variables
is fixed in advance. Let
\[
 F\in\Z[Y_1,\ldots,Y_m,T_1,\ldots,T_k]^e,\qquad x=(y_0,t_0),
\]
and assume $F(x)=q^Lr_x$ with $q\ge1$ and $L\ge1$. Write the actual
reduced Jacobian as
\[
 DF(x)\bmod q=(J_x\ K_x),\qquad
 J_x:S_q^m\to S_q^e,\quad K_x:S_q^k\to S_q^e.
\]
For a prescribed parameter digit $v\in S_q^k$, let $\mathcal L_{x,L,v}$
be the set of next residue roots over $x$ whose parameter coordinates are
$t_0+q^L\widetilde v$ modulo $q^{L+1}$.

\begin{theorem}[Prescribed-parameter obstruction and count]
\label{thm:prescribed-obstruction}
Define
\[
 o_x(v)=[\overline{r_x}+K_xv]\in S_q^e/\im J_x.
\]
Then $\mathcal L_{x,L,v}$ is nonempty if and only if $o_x(v)=0$.
When nonempty, it is a torsor for $\ker J_x$. If $q=p$ is prime and
$\rho_x=\rank_{\F_p}J_x$, then
\begin{equation}\label{eq:prescribed-count}
 \#\mathcal L_{x,L,v}
 =\begin{cases}p^{m-\rho_x},&o_x(v)=0,\\0,&o_x(v)\ne0,\end{cases}
 \qquad \#(\F_p^e/\im J_x)=p^{e-\rho_x}.
\end{equation}
The criterion holds uniformly for every prescribed $v$ without a
surjectivity hypothesis.
\end{theorem}

\begin{proof}
The actual next point has the form
$x+q^L(\widetilde u,\widetilde v)$, with solved digit $u\in S_q^m$.
Its Taylor correction equation is
\[
 \overline{r_x}+J_xu+K_xv=0.
\]
Solvability is exactly vanishing of the displayed cokernel class. Restrict
the residue-digit equivalence to the prescribed parameter coordinates;
this identifies the actual lift fiber with the solved-digit fiber.
Translation by $\ker J_x$ is free and transitive on any nonempty such
fiber. The prime-field count follows from rank-nullity as before.
\end{proof}

The lower-root-only statements, including construction of the residual,
are \lean{Komorebi.PolynomialPrescribedObstruction.finite_prescribed_obstruction}
and \lean{Komorebi.PolynomialPrescribedObstruction.finite_prescribed_lift_count}.
For $F(Y,T)=T$ at the dyadic origin, the full Jacobian $(0\ 1)$ is
surjective but $J_x=0$. Prescribing $v=1$ forces $T=2\pmod4$, so no root
exists; prescribing $v=0$ gives two choices of $Y$ modulo four over the
initial zero residue. This checks why the solved block, rather than merely
the full Jacobian, appears in the criterion.

\begin{theorem}[Every compatible integer parameter path]
\label{thm:prescribed-integer-tower}
In the preceding notation, assume $J_x$ is surjective. For every supplied
integer path $t_n\in\Z^k$ satisfying
\[
 t_0=\operatorname{free}(x),\qquad
 t_{n+1}-t_n\in q^{L+n}\Z^k,
\]
there is a compatible polynomial tower $z_n\in\Z^{m+k}$ with
\[
 z_0=x,\quad \operatorname{free}(z_n)=t_n,\quad
 F(z_n)\in q^{L+n}\Z^e,\quad
 z_{n+1}-z_n\in q^{L+n}\Z^{m+k}.
\]
The free-coordinate equality is exact in integers. Both reduced Jacobian
blocks remain fixed, and $z_b-z_a\in q^{L+a}\Z^{m+k}$ whenever $a\le b$.
\end{theorem}

\begin{proof}
At level $N=L+n$, write the supplied parameter increment exactly as
$t_{n+1}-t_n=q^Nw$ with $w\in\Z^k$. Surjectivity of the solved block
gives a digit $u$ satisfying
$J_{z_n}u=-\overline{F(z_n)/q^N}-K_{z_n}\overline w$.
Choose an integer representative $\widetilde u$ and set
$z_{n+1}=z_n+q^N(\widetilde u,w)$. Its parameter coordinates equal the
supplied target exactly. Taylor expansion proves the next root condition,
while congruence modulo $q$ preserves the full reduced Jacobian.
Dependent recursion constructs the sequence with the exact parameter
coordinate and the fixed solved Jacobian in its state. The earlier
cross-level argument proves all retained integer precisions.
\end{proof}

This is
\lean{Komorebi.PolynomialPrescribedIntegerTower.exists_prescribed_tower_strong}.
In particular, no assumption about solvability at each later level has been
added to the initial surjectivity condition.

\begin{theorem}[Arbitrary $p$-adic parameter endpoints]
\label{thm:prescribed-padic-tower}
Let $q=p$ be prime in Theorem~\ref{thm:prescribed-integer-tower}.
For every $\tau\in t_0+p^L\Zp^k$, there are an integer tower $z_n$ and
$b\in\Zp^{m+k}$ such that
\[
 F(b)=0,\qquad\operatorname{free}(b)=\tau,\qquad
 z_n\longrightarrow b,\qquad b-z_n\in p^{L+n}\Zp^{m+k}.
\]
The initial value is $z_0=x$. At every stage and for every proposed
parameter digit $v\in\F_p^k$, the corresponding next residue fiber has
exactly $p^{m-e}$ elements. If $m=e$, then for each $\tau$ there is exactly
one such root in $B_L(x)$, independently of all choices of integer
approximants and correcting digits.
\end{theorem}

\begin{proof}
Set the first integer parameter approximation equal to the prescribed
$t_0$. For $n\ge1$, use canonical representatives of
$\tau\bmod p^{L+n}$. The initial congruence and truncation precision give
\[
 \tau-t_n\in p^{L+n}\Zp^k.
\]
Subtracting consecutive relations gives compatibility in $\Zp$.
Faithfulness of integer divisibility under $\Z\hookrightarrow\Zp$
gives compatibility in $\Z$, and the same norm bounds give $t_n\to\tau$.
Theorem~\ref{thm:prescribed-integer-tower} and completion now produce $b$
with exact polynomial vanishing and all retained precisions. Each free
coordinate has both the corresponding coordinate of $b$ and that of
$\tau$ as its limit, proving exact parameter equality.

The solved Jacobian remains surjective, so every shifted obstruction
$[\overline r+K_{z_n}v]$ vanishes. Formula~\eqref{eq:prescribed-count}
therefore gives $p^{m-e}$ choices for every parameter digit at every stage.
When $m=e$, the fixed-coordinate chart in the proof of
Theorem~\ref{thm:continuous-parameter-chart} contains every root in the
initial ball with the specified parameters. Its uniqueness proves the
endpoint assertion. This is uniqueness of the actual endpoint, not of the
integer representatives used to approach it.
\end{proof}

The complete interfaces are
\lean{Komorebi.PolynomialPrescribedPadicTower.exists_prescribed_padic_root_with_tower}
and \lean{Komorebi.PolynomialPrescribedPadicTower.exists_unique_prescribed_padic_root}.
The separate supplied-path theorem
\lean{Komorebi.PolynomialPrescribedPadicTower.exists_root_with_prescribed_path}
accepts any compatible integer parameter path with limit $\tau$.
For $m>e$, endpoint uniqueness is false in general: the system
$F(Y_1,Y_2,T)=Y_1$ has two distinct roots $(0,0,0)$ and $(0,2,0)$ in the
same dyadic ball with identical parameter value.

\subsection{Weighted lifting for a nonunit Jacobian}

The remaining theorems in this section allow arbitrary coefficients and
centers in $R=\Zp$. They concern square solved Jacobians invertible over
$\Qp$, even when their reductions modulo $p$ are singular. The exponent
controlling an integral inverse is retained explicitly.

\begin{theorem}[Weighted nonlinear lifting]\label{thm:weighted-lifting}
Let $p$ be prime, $d,s,r\in\N$, $F\in R[X_1,\ldots,X_d]^d$, and
$a\in R^d$. Set $J=DF(a)$. Suppose
\begin{equation}\label{eq:weighted-hypotheses}
 J\in\operatorname{GL}_d(\Qp),\qquad
 p^sJ^{-1}\in M_d(R),\qquad F(a)\in p^rR^d,\qquad r>2s.
\end{equation}
Then there is exactly one root $b$ of $F$ in $a+p^{s+1}R^d$, and it
satisfies
\begin{equation}\label{eq:weighted-displacement}
 b-a\in p^{r-s}R^d.
\end{equation}
Uniqueness holds throughout the full ball of depth $s+1$.
\end{theorem}

\begin{proof}
Let $A=p^sJ^{-1}\in M_d(R)$. Inclusion into $\Qp$ and its injectivity
give the two integral identities
\[
 AJ=JA=p^sI_d.
\]
These identities are also an equivalent integral certificate for the
first two hypotheses of \eqref{eq:weighted-hypotheses}. In particular,
$A$ is injective: $Aw=0$ implies $p^sw=JAw=0$, hence $w=0$.

Precondition the actual polynomial list by $A$ and set $c=p^{s+1}$.
Its derivative at $a$ is $AJ=p^sI_d$. Lemma~\ref{lem:polynomial-taylor}
gives a polynomial vector $Q$ with $Q(0)=0$ and
\[
 (AF)(a+cu)=(AF)(a)+p^{2s+1}u+p^{2s+2}Q(u).
\]
Since $A$ is integral, $(AF)(a)\in p^rR^d$. The strict integer inequality
$r>2s$ permits $k=r-(2s+1)\ge0$ and a vector $C\in p^kR^d$ such that
$(AF)(a)=p^{2s+1}C$. The normalized root equation is therefore
\begin{equation}\label{eq:weighted-fixed-point}
 u=-C-pQ(u).
\end{equation}

Integral polynomial evaluation is $1$-Lipschitz for the supremum
$p$-adic metric: a common divisibility of input differences is preserved
under addition and multiplication by integral elements. Consequently the
right side of \eqref{eq:weighted-fixed-point} has contraction constant
$|p|_p<1$ on the complete space $R^d$. It has a unique fixed point $u$.

To obtain the sharper displacement, begin with $u\in R^d$. If
$u\in p^nR^d$, then $Q(0)=0$ and polynomial congruence give
$Q(u)\in p^nR^d$. For $n<k$, both $C$ and $pQ(u)$ are divisible by
$p^{n+1}$. Equation~\eqref{eq:weighted-fixed-point} therefore improves
the divisibility of $u$ by one digit. Induction gives $u\in p^kR^d$.
Thus $b=a+p^{s+1}u$ satisfies \eqref{eq:weighted-displacement}, because
$s+1+k=r-s$. Its normalized preconditioned value vanishes, and injectivity
of $A$ gives the actual equation $F(b)=0$.

Every competing root $z\in a+p^{s+1}R^d$ can be written
$z=a+p^{s+1}w$. Its actual root equation implies the same fixed-point
equation after cancellation of the nonzero scalar $p^{2s+1}$.
Contraction uniqueness gives $w=u$, proving uniqueness on the entire
stated ball.
\end{proof}

The main entry points are
\lean{Komorebi.PolynomialWeightedLift.exists_weighted_root} and
\lean{Komorebi.PolynomialWeightedLift.exists_unique_weighted_root}.
The inverse hypothesis in Lean specifies an integral matrix whose image
is exactly $p^sJ^{-1}$; it does not substitute a potentially larger
determinant valuation for $s$.

\begin{example}[Strict threshold and displacement boundary]
\label{ex:weighted-boundary}
For $p=2$, $F(X)=X^2+3$, and $a=1$, one has $F(a)=4$ and $DF(a)=2$.
Thus $s=1$ and $r=2$ satisfy every displayed hypothesis except the strict
inequality. No root exists in $\Z_2$, since none exists modulo eight.
Hence $r\ge2s$ cannot replace $r>2s$ uniformly. Also,
$F(X)=p^sX+p^r$ has root $-p^{r-s}$, whose displacement is not divisible
by $p^{r-s+1}$. The guaranteed displacement cannot uniformly gain a digit.
\end{example}

\subsection{First-order coefficient and parameter motion}

\begin{theorem}[Weighted coefficient deformation]\label{thm:coefficient-motion}
Let $F,H\in R[X_1,\ldots,X_d]^d$, $F(a)=0$, and $J=DF(a)$. Suppose
$J\in\operatorname{GL}_d(\Qp)$, $p^sJ^{-1}\in M_d(R)$, and $M>2s$.
For $G=F+p^MH$, there is exactly one root $b\in a+p^{s+1}R^d$.
It satisfies $b-a\in p^{M-s}R^d$ and, in $\Qp^d$,
\begin{equation}\label{eq:coefficient-first-order}
 b-a=-p^MJ^{-1}H(a)+\varepsilon,
 \qquad\varepsilon\in p^{2M-3s}R^d.
\end{equation}
Only the original Jacobian $DF(a)$ is subject to an inverse hypothesis.
\end{theorem}

\begin{proof}
Retain the original certificate $A=p^sJ^{-1}$ and put $c=p^{s+1}$.
Taylor expansion and a polynomial difference identity yield polynomial
vectors $Q_F,Q_H$, both zero at the origin, such that
\[
 (AF)(a+cu)=p^{2s+1}u+p^{2s+2}Q_F(u),\qquad
 (AH)(a+cu)=AH(a)+cQ_H(u).
\]
Consequently
\[
 (AG)(a+cu)=p^{2s+1}\bigl(C+u+pQ(u)\bigr),
\]
where
\[
 C=p^{M-(2s+1)}AH(a),\qquad
 Q=Q_F+p^{M-s-1}Q_H,\qquad Q(0)=0.
\]
All exponents are nonnegative. The same contraction and divisibility
argument as in Theorem~\ref{thm:weighted-lifting} gives the actual root,
the displacement $p^{M-s}$, and full-ball uniqueness. This constructs the
root without introducing an inverse hypothesis for $DG(a)$.

For the error estimate, let $\delta=b-a$ and $D=M-s$. Taylor's theorem
for the actual preconditioned polynomial gives
\[
 E_F=(AF)(b)-(AF)(a)-p^s\delta\in p^{2D}R^d.
\]
Polynomial congruence gives
$E_H=(AH)(b)-(AH)(a)\in p^DR^d$.
The equations $F(a)=0$ and $G(b)=0$ imply
\[
 p^s\delta+p^MAH(a)=-E_F-p^ME_H\in p^{2D}R^d,
\]
because $M+D\ge2D$. Cancel $p^s$ in the integral domain $R$ to obtain
\[
 \delta+p^{M-s}AH(a)\in p^{2M-3s}R^d.
\]
This integral vector is $\varepsilon$. After inclusion into $\Qp$,
$p^{M-s}A=p^MJ^{-1}$ gives \eqref{eq:coefficient-first-order}.
\end{proof}

The complete theorem is
\lean{Komorebi.PolynomialRootMotion.exists_unique_root_motion}.
The error depth exceeds the displacement depth by $M-2s\ge1$. It is a
guaranteed divisibility, not an assertion of exact valuation for every
system.

\begin{theorem}[Simultaneous coefficient and parameter deformation]
\label{thm:joint-deformation}
Let $F,H\in R[Y_1,\ldots,Y_e,T_1,\ldots,T_k]^e$ and
$a=(y_0,t_0)\in R^{e+k}$ satisfy $F(a)=0$. Define the actual blocks
\[
 J=D_YF(a)\in M_e(R),\qquad K=D_TF(a)\in M_{e\times k}(R).
\]
Suppose $J\in\operatorname{GL}_e(\Qp)$,
$p^sJ^{-1}\in M_e(R)$, and $M>2s$. For every $v\in R^k$, set
\[
 \tau=t_0+p^Mv,\qquad G=F+p^MH.
\]
There is exactly one $y\in y_0+p^{s+1}R^e$ with $G(y,\tau)=0$.
It satisfies
\begin{equation}\label{eq:joint-first-order}
 \begin{aligned}
  y-y_0&\in p^{M-s}R^e,\\
  y-y_0&=-p^MJ^{-1}\bigl(H(a)+Kv\bigr)+\varepsilon,
  \qquad\varepsilon\in p^{2M-3s}R^e.
 \end{aligned}
\end{equation}
The equality involving $J^{-1}$ is in $\Qp^e$, with an integral error
vector. The parameter equality is exact.
\end{theorem}

\begin{proof}
Use the actual augmented polynomial systems
\[
 \mathcal F(Y,T)=(F(Y,T),T-t_0),\qquad
 \mathcal H(Y,T)=(H(Y,T),-v).
\]
Then $\mathcal F(a)=0$, and its actual derivative is
\[
 B=D\mathcal F(a)=\begin{pmatrix}J&K\\0&I_k\end{pmatrix}.
\]
For $A=p^sJ^{-1}\in M_e(R)$ and $c=p^s$, define
\[
 C=\begin{pmatrix}A&-AK\\0&cI_k\end{pmatrix}.
\]
The identities $AJ=JA=cI_e$ give $CB=BC=cI_{e+k}$. In the second
product, the upper right block is $-JAK+cK=0$; in the first, it is
$AK-AK=0$. Thus augmentation preserves the original inverse bound $s$.

Theorem~\ref{thm:coefficient-motion} gives a unique root $b$ of
$\mathcal F+p^M\mathcal H$ in $a+p^{s+1}R^{e+k}$, with displacement
depth $M-s$. Its last equations say exactly
$b_T=t_0+p^Mv=\tau$. Its first equations are $G(b)=0$.
The integral first-order expression uses
\[
 C\binom{H(a)}{-v}
 =\binom{A(H(a)+Kv)}{-p^sv}.
\]
Projection onto the solved coordinates gives
\[
 b_Y-y_0+p^{M-s}A(H(a)+Kv)\in p^{2M-3s}R^e.
\]
The same inclusion identity as before converts this to
\eqref{eq:joint-first-order}.

Finally, if $z\in y_0+p^{s+1}R^e$ is any competing solved root with
parameters $\tau$, then $(z,\tau)$ belongs to the full augmented ball.
Indeed, $M>2s$ implies $M\ge s+1$. Full-point uniqueness therefore gives
$(z,\tau)=b$. This comparison proves uniqueness on exactly the solved
ball stated in the theorem.
\end{proof}

The exact solved-coordinate statement is
\lean{Komorebi.PolynomialJointDeformation.exists_joint_deformation},
with its unique-existence interface
\lean{Komorebi.PolynomialJointDeformation.exists_unique_joint_deformation}.
The triangular inverse identities are checked over arbitrary commutative
rings before their $p$-adic specialization.

\begin{example}[First-order sign and error precision]\label{ex:motion-boundary}
For $p=2$, $s=0$, $M=2$, $F(Y,T)=Y+T$, $H=1$, $a=0$, and $v=1$,
the prescribed parameter is $\tau=4$ and the solved root is $y=-8$.
Here $J=K=1$, and the forcing term is $H(a)+Kv=2$. Omitting the parameter
term would predict $-4$, and reversing its sign would predict $0$.
The corresponding errors $-4$ and $-8$ are both outside
$2^{2M}\Z_2=16\Z_2$. These are explicit checked sign controls.

The error exponent cannot uniformly be increased. For
$F(Y,T)=Y+Y^2+T$, $H=0$, $a=0$, $p=2$, $s=0$, $M=1$, and $v=1$,
there is an even root $y$ of $y+y^2+2=0$. Its equation forces
$y\equiv2\pmod4$ and $y+2=-y^2$, so
\[
 4\mid y+2,\qquad 8\nmid y+2.
\]
Thus the error depth $2M-3s=2$ is attained even for parameter-only motion.
The same quadratic with a constant coefficient perturbation gives the
coefficient-only boundary. These examples refute a uniform extra error
digit; they do not assert optimality for every individual choice of
$p,s,M$.
\end{example}

The combined results establish finite singular obstructions, smooth
compatible towers, full continuous root charts after a suitable coordinate
choice, and quantitative deformation for a rationally invertible solved
Jacobian. Rationally rank-deficient persistence remains a separate problem.
The analytic solution-map theorem in the next section supplies the
regularity conclusion with the same actual derivative blocks.

\section{Analytic regularity of the local solution map}
\label{sec:regularity}
The deformation estimates quantify the displacement of an exact root under
small changes. We now identify the local parameter dependence as analytic.
The construction uses the actual polynomial evaluation map and its actual
Jacobian. Its domain is a vector space over the fraction field. In particular,
$\Zp$ is not used as a vector space over $\Qp$.

Let $\mathbb K$ be a complete nontrivially normed field. Let $e,k\in\N$, and
write the variables as $(Y,T)\in\mathbb K^e\times\mathbb K^k$. For
\[
 F\in\mathbb K[Y_1,\ldots,Y_e,T_1,\ldots,T_k]^e,
 \qquad a=(y_0,t_0),
\]
define the matrices of evaluated formal partial derivatives
\begin{equation}
 J=D_YF(a)\in M_e(\mathbb K),\qquad
 K=D_TF(a)\in M_{e\times k}(\mathbb K).
 \label{eq:regularity-blocks}
\end{equation}
Rows index equations and columns index variables. We first establish that
these formal derivatives are the derivatives of the evaluation map.

\begin{lemma}[Actual polynomial derivative]
\label{lem:actual-polynomial-derivative}
For $P\in\mathbb K[X_1,\ldots,X_n]$, polynomial evaluation is analytic at
every point. At $a\in\mathbb K^n$ its strict Fr\'echet derivative is
\[
 v\longmapsto\sum_{j=1}^n (\partial_jP)(a)v_j.
\]
For a finite system, its strict derivative is the continuous linear map
represented by the matrix of evaluated formal partial derivatives.
\end{lemma}
\begin{proof}
For the derivative statement, use polynomial induction on constants,
addition, and multiplication by a variable. Constants have zero derivative.
The sum rule agrees with the additivity of formal partial differentiation.
For $PX_j$, the product rule gives
\[
 d(PX_j)_a(v)=a_j\,dP_a(v)+P(a)v_j,
\]
which agrees with the evaluated formal product rule for every coordinate.
The finite product rule for maps to coordinate spaces then assembles the
system derivative. Each matrix defines a continuous linear map because the
coordinate spaces are finite-dimensional over the complete field.
Polynomial analyticity follows from the corresponding polynomial induction
in Mathlib and its finite-coordinate assembly.
\end{proof}
The derivative bridge and analytic statements are
\lean{Komorebi.PolynomialAnalyticEvaluation.hasStrictFDerivAt_evalSystem},
\lean{Komorebi.PolynomialAnalyticEvaluation.analyticAt_evalSystem}, and
\lean{Komorebi.PolynomialAnalyticEvaluation.analyticOnNhd_evalSystem}.

\begin{lemma}[The augmented derivative]
\label{lem:augmented-analytic-inverse}
If $J\in\operatorname{GL}_e(\mathbb K)$, the actual polynomial map
\[
 \Phi(Y,T)=(F(Y,T),T)
\]
has invertible strict derivative at $a$, with
\begin{equation}
 D\Phi(a)=\begin{pmatrix}J&K\\0&I_k\end{pmatrix},\qquad
 D\Phi(a)^{-1}=\begin{pmatrix}J^{-1}&-J^{-1}K\\0&I_k\end{pmatrix}.
 \label{eq:analytic-block-inverse}
\end{equation}
The local inverse supplied by this derivative is analytic near $\Phi(a)$.
\end{lemma}
\begin{proof}
The lower equations are the actual parameter variables. Differentiating
them gives the blocks $0$ and $I_k$. Lemma~\ref{lem:actual-polynomial-derivative}
identifies the remaining blocks with~\eqref{eq:regularity-blocks}.
Multiplication of the two matrices in~\eqref{eq:analytic-block-inverse},
in both orders, gives the identity. The resulting linear equivalence is
continuous, as is its inverse. Apply the strict inverse-function theorem to
$\Phi$. The analytic inverse theorem applies to the same open partial
homeomorphism, because $\Phi$ is analytic and its derivative is that
continuous linear equivalence. It supplies a convergent local power-series
representation of this inverse.
\end{proof}
The actual augmented Jacobian is checked by
\lean{Komorebi.PolynomialJointSystem.jacobian_augmented}.
The continuous inverse block and its restriction are
\lean{Komorebi.PolynomialAnalyticBlocks.upperEquiv} and
\lean{Komorebi.PolynomialAnalyticBlocks.inverse_solved_zero}.
Analyticity of the same constructed inverse is
\lean{Komorebi.PolynomialAnalyticInverse.analyticAt_inverse}.

\begin{theorem}[Analytic local root graph]
\label{thm:analytic-local-graph}
Suppose $F(a)=0$ and $J\in\operatorname{GL}_e(\mathbb K)$. There are open
neighborhoods $t_0\in U\subseteq\mathbb K^k$ and
$y_0\in V\subseteq\mathbb K^e$, and a map
$\psi:\mathbb K^k\to\mathbb K^e$, such that
\begin{gather}
 \psi(t_0)=y_0,\qquad \psi(U)\subseteq V,\qquad
 \psi\text{ is analytic at every point of }U,\label{eq:analytic-psi}\\
 F(y,t)=0\quad\Longleftrightarrow\quad y=\psi(t)
 \qquad(t\in U,\ y\in V).\label{eq:analytic-root-graph}
\end{gather}
At the base parameter, the strict Fr\'echet derivative is
\begin{equation}
 \boxed{D\psi(t_0)=-J^{-1}K.}
 \label{eq:implicit-derivative}
\end{equation}
In particular, $\psi$ is differentiable on $U$.
\end{theorem}
\begin{proof}
Let $g$ be the analytic local inverse of $\Phi$ supplied by
Lemma~\ref{lem:augmented-analytic-inverse}. Since $F(a)=0$,
$\Phi(a)=(0,t_0)$ and $g(0,t_0)=a$. Define
\[
 \psi(t)=\pi_Yg(0,t),
\]
where $\pi_Y$ projects to the solved coordinates. This gives the base value
and analyticity at $t_0$. The chain rule and
\eqref{eq:analytic-block-inverse} give
\[
 D\psi(t_0)v
 =\pi_Y\begin{pmatrix}J^{-1}&-J^{-1}K\\0&I_k\end{pmatrix}
       \binom{0}{v}
 =-J^{-1}Kv.
\]

For completeness, the neighborhood construction must provide an entire
product graph. The source of the open partial homeomorphism defined by
$\Phi$ is open around $a$; its inverse is $g$. Choose open sets $V$ and $U_0$
containing $y_0$ and $t_0$ such that
$V\times U_0$ lies in that source. Restrict $U_0$ further by the preimage of
an open analytic inverse neighborhood under $t\mapsto(0,t)$, the open locus
where this same $\psi$ is analytic, and the interior of $\psi^{-1}(V)$.
All of these sets contain $t_0$. For the last set, use the continuity of
$\psi$ at $t_0$ and $\psi(t_0)=y_0\in V$. Their intersection is an open
neighborhood $U$ with $\psi(U)\subseteq V$.

If $t\in U$, $y\in V$ and $F(y,t)=0$, then
$\Phi(y,t)=(0,t)$. The left inverse identity forces
$g(0,t)=(y,t)$, hence $y=\psi(t)$. Conversely, the right inverse identity
$\Phi(g(0,t))=(0,t)$ says both that its solved equations vanish and that its
parameter coordinates are exactly $t$. Thus
$g(0,t)=(\psi(t),t)$ and $F(\psi(t),t)=0$.
Analyticity at every point of $U$ implies differentiability there.
\end{proof}
This is
\lean{Komorebi.PolynomialAnalyticImplicit.exists_analytic_implicit}.
The expanded statement contract also includes differentiability explicitly:
\lean{Komorebi.PolynomialRegularityContracts.actual_padic_analytic_root_graph}.

\begin{theorem}[Regularity of a supplied root selection]
\label{thm:selection-regularity}
Under the hypotheses of Theorem~\ref{thm:analytic-local-graph}, let
$h:\mathbb K^k\to\mathbb K^e$ be continuous at $t_0$, with
$h(t_0)=y_0$. Suppose $F(h(t),t)=0$ for every $t$ in some neighborhood of
$t_0$. Then $h$ and $\psi$ agree in a neighborhood of $t_0$. Consequently
$h$ is analytic there as a germ and has strict derivative
$Dh(t_0)=-J^{-1}K$.
\end{theorem}
\begin{proof}
Continuity at $t_0$ places $h(t)$ inside $V$ for all sufficiently close
parameters. Intersect this neighborhood with $U$ and the neighborhood on
which the root equation holds. Equation~\eqref{eq:analytic-root-graph} gives
$h(t)=\psi(t)$ throughout the intersection. Analyticity and strict
differentiability at a point are unchanged by agreement on a neighborhood.
\end{proof}
The formal result is
\lean{Komorebi.PolynomialAnalyticSelection.continuous_selection_regularity}.
Its formulation even derives $F(a)=0$ from the eventual root equation and
the specified base value.

\begin{corollary}[Regularity throughout an open parameter set]
\label{cor:selection-open-regularity}
Let $U\subseteq\mathbb K^k$ be open, and let $h$ be continuous on $U$.
Suppose that for every $t\in U$,
\[
 F(h(t),t)=0,
 \qquad J_t=D_YF(h(t),t)\in\operatorname{GL}_e(\mathbb K).
\]
Then $h$ is analytic at every point of $U$ and its strict derivative is
\[
 Dh(t)=-J_t^{-1}D_TF(h(t),t)\qquad(t\in U).
\]
\end{corollary}
\begin{proof}
At each $t\in U$, openness turns continuity on $U$ into continuity at $t$
and makes the root equation hold on a neighborhood of $t$. Apply
Theorem~\ref{thm:selection-regularity} at $(h(t),t)$.
\end{proof}
This is
\lean{Komorebi.PolynomialAnalyticSelection.continuous_selection_analyticOnNhd}.
The pointwise invertibility assumption is part of this corollary.

\begin{corollary}[Integral coefficients and centers]
\label{cor:integral-analytic-graph}
Let $p$ be any prime. Suppose $F\in\Zp[Y,T]^e$, $a\in\Zp^{e+k}$,
and $F(a)=0$. Let $\iota:\Zp\hookrightarrow\Qp$ be the inclusion.
If $\iota(D_YF(a))\in\operatorname{GL}_e(\Qp)$, the polynomial system
obtained by including the original coefficients in $\Qp$ has the analytic
local root graph of Theorem~\ref{thm:analytic-local-graph} at $\iota(a)$.
The strict derivative at the base is
\[
 -\iota(D_YF(a))^{-1}\iota(D_TF(a)).
\]
The open neighborhoods and the output map are in the ambient spaces over
$\Qp$.
\end{corollary}
\begin{proof}
Coefficient inclusion commutes with evaluation and formal partial
differentiation. Consequently both Jacobian blocks after inclusion are the
entrywise inclusions of the original blocks, and the original exact root
remains an exact root. Apply Theorem~\ref{thm:analytic-local-graph} to these
actual included polynomials.
\end{proof}
This is
\lean{Komorebi.PolynomialRegularityIntegral.exists_analytic_integral_implicit}.
It requires no invertibility modulo $p$. In particular, it includes solved
Jacobians singular modulo $p$ but invertible over the fraction field.

\begin{example}[The derivative sign and a nonlinear control]
For $F(Y,T)=Y+T$, the exact branch is $\psi(T)=-T$ and $J=K=1$.
The derivative is $-1$. Uniqueness of the derivative rules out $+1$ over
$\Qp$, including $p=2$. The nonlinear polynomial $F(Y,T)=Y+Y^2+T$
has $F(0,0)=0$ and $J=K=1$ there. The actual formal theorem is instantiated
on this polynomial to construct an analytic local root graph of derivative
$-1$; a branch is not supplied as an assumption.
The declarations and the deliberately rejected positive-sign statement are
in the regularity control and statement-contract files.
\end{example}

The results cover arbitrary finite $e$ and $k$, including zero. The local
inverse machinery explicitly handles a zero-dimensional domain. These
analytic neighborhoods have no new uniform numerical radius. The earlier
weighted lifting and joint deformation theorems retain their quantitative
radii and error bounds, while the present statements establish regularity
of the local solution germ. A direct Lean identification with each
historical subtype-valued chart is an additional interface task. General
persistence at a fraction-field singular solved block is a separate problem.

\section{Remaining questions and the boundary of the results}
\label{sec:open}

The external Nishibuchi conjecture, in its stated parameter range, is resolved
by the arithmetic results above. The following questions concern extensions,
sharper classifications, and the original native seed-calculus program. They
are not missing steps in that resolved conjecture. The local analyticity and
differentiability theorem is also complete under its stated invertibility
hypothesis.

\subsection{Singular persistence beyond an invertible solved block}
\label{subsec:singular-open}

The weighted lifting theorem allows the solved Jacobian to be noninvertible
modulo $p$ provided that it is invertible over $\Qp$ with a controlled inverse.
This includes a substantial class of lifts often called singular modulo $p$.
The remaining problem concerns a different situation: the Jacobian is
rank-deficient even over $\Qp$, and the current inverse construction cannot
apply.

For a fixed polynomial system $F\in\Z[X_1,\ldots,X_d]^e$, let
\[
\mathcal L_{F,L}(\bar x)
=\{x'\in(\Z/p^{L+1}\Z)^d:
F(x')=0,\ x'\equiv\bar x\pmod{p^L}\}.
\]
The finite obstruction theorem already decides whether this set is empty
from the actual residual digit and reduced Jacobian. If it is nonempty, its
cardinality and affine kernel structure are determined. A vanishing class
at one level says only that this one step is possible.

For example, $X^2-5$ has the root $1$ modulo $2$, and both $1$ and $3$ lift it
to roots modulo $4$. Every odd square is $1$ modulo $8$, so neither branch
continues to a root modulo $8$. This direct arithmetic example illustrates
why a successful first correction is not a persistence theorem.

The open task is to find useful structural hypotheses and higher obstruction
data that control an entire singular branch, including its behavior under
coefficient and parameter deformation. One must establish those conditions
at successive levels, or establish a suitable all-level existence statement
from which compatible branches follow. The finite obstruction calculus gives
a way to test each level. It does not supply an unconditional singular
persistence result. Multiple branches, later obstruction, and changes in the
rank of the Jacobian must remain possible in any proposed general statement.

\subsection{Native 13.2: what presentation invariance must mean}
\label{subsec:native132}

The original conjecture asks for a canonical finite obstruction packet whose
vanishing is necessary and sufficient for a one-level lift in a finitely
presented arithmetic seed system. For a fixed finite polynomial presentation,
the paper now supplies the precise obstruction
\[
o_{F,\bar x}=[r]\in
\coker\!\left(J_{F,\bar x}:\F_p^d\longrightarrow\F_p^e\right),
\]
where $r$ is the residual digit. Its independence from the chosen integer
representative is proved. That independence is different from invariance
under a change of equations or presentation.

The distinction cannot be removed by declaring the raw cokernel canonical.
The presentations $F=(X)$ and $G=(X,0)$ describe the same zero locus over every
coefficient ring. At zero their Jacobians are, respectively,
\[
\F_p\xrightarrow{\,1\,}\F_p,
\qquad
\F_p\xrightarrow{\,(1,0)\,}\F_p^2.
\]
Their raw cokernels are $0$ and $\F_p$. Both distinguished obstruction classes
vanish, but the two ambient obstruction groups are not isomorphic. This is
an elementary presentation counterexample, not a refutation of every possible
intrinsic obstruction construction.

A complete version of Native 13.2 therefore needs the following data before
its full statement can be audited:
\begin{enumerate}
\item a category of seed objects, including how the operators, relations,
projectors, weights, and topology enter an object;
\item specified changes of presentation and a proof that they preserve the
actual lift problem;
\item an obstruction object that respects dependencies among the relations,
together with comparison maps or a universal property;
\item a precise meaning of ``finite'' and ``canonical'', and a proof that
vanishing is equivalent to the intended one-step continuation.
\end{enumerate}
For commutative algebra presentations, retaining the relation module rather
than all listed equations is a natural route. Extending such a construction
to the entire seed vocabulary requires an encoding and a proof that lifts
are preserved. The consolidation does not claim that missing implication.

\subsection{Native 13.3: finite compatibility, completion, and descent}
\label{subsec:native133}

The original proposed form is
\[
X\simeq X_2^{\mathrm{wt}}\times_C\prod_{p>2}'X_p.
\]
It leaves $X$, the distinguished subsets defining the restricted product,
the maps to $C$, the ambient category, and the universal meaning of
``smallest'' unspecified. In particular, one must decide whether $X$ is a
global object or a completed object.

There is now an exact linear compatibility result. If two local observations
$J:U\to V$ and $K:U\to W$ are linear maps from the \emph{same} correction
module, define
\[
C_{J,K}=U/(\ker J+\ker K).
\]
The attainable observations $\im J$ and $\im K$ have canonical labels in this
quotient. Equality of labels is equivalent to simultaneous realization by a
single correction. The joint image is exactly the fiber product over
$C_{J,K}$, with the proved linear universal property.
The declarations are
\begingroup\small\raggedright
\lean{Komorebi.SharedCorrection.labels_eq_iff_common},
\lean{Komorebi.SharedCorrection.joint_range_eq_fiber}, and
\lean{Komorebi.SharedCorrection.common_universal}.
\par\endgroup
This result assumes a common correction module and actual observation maps.
An arithmetic application must construct them and prove that their notion
of compatibility is the one required by the application.

\paragraph{Why the global/completed distinction is necessary.}
Consider constant sequences with integer initial value. Their global parameter
space is $\Z$, while their local parameter spaces are $\Z_p$. Under the
canonical local maps, the image of the global objects is the diagonal
\[
D=\{(n)_p:n\in\Z\}\subseteq\prod_p\Z_p.
\]
This diagonal is dense: a basic open set specifies finitely many congruences,
and the Chinese remainder theorem provides an integer satisfying them. It is
proper: the tuple with $2$-adic coordinate zero and all odd-adic coordinates
one is not the image of an integer. Indeed an integer zero in $\Z_2$ must be
zero, contradicting its $3$-adic coordinate.

Consequently this canonical global diagonal cannot be an equality locus
\[
\{(a,b):f(a)=g(b)\}
\subseteq \Z_2\times\prod_{p>2}\Z_p
\]
for continuous maps to a Hausdorff space $C$. Such an equality locus is closed,
as the inverse image of the closed diagonal in $C\times C$, whereas $D$ is
proper and dense. This elementary topological argument is retained from the
native progress note as a prose argument. It is not a new Lean-certified
result, and it concerns exactly the stated topological interpretation.

Finite CRT reconstruction and the familiar completed object
$\widehat\Z=\prod_p\Z_p$ address a different question. They do not recover
which local tuples arise from a single global integer. A complete form of
Native 13.3 must specify finite-level compatibility, transition maps and
completion, and any additional global descent condition. Recurrence,
spectral, or Frobenius information must appear as actual maps or constraints,
not only as a proposed description of $C$.

\subsection{Arithmetic refinements}
\label{subsec:arithmetic-open}

Write $U_{N,j}(n)=j!C_{N,j}(n)$ and $U_j=U_{4,j}$, and use
$v_p(0)=\infty$ when discussing valuation bounds. The following classification
questions remain open in this program.

\begin{enumerate}
\item \textbf{Higher dyadic differences.}
Determine the exact value of
\[
v_2\!\left(\Delta_u^kU_j(n)\right),\qquad k\geq2,
\]
with explicit treatment of zero differences. The all-order lower bounds and
the exact first-displacement theorem do not determine this value: subtraction
of two first differences can introduce further cancellation. The first
concrete target is order $2$ in one fixed address.

\item \textbf{Least periods beyond $N=4$.}
For further packet families, determine the least positive $P$ for which
$U_{N,j}(n+P)\equiv U_{N,j}(n)\pmod{p^r}$ for every $n$ in the claimed range.
The general sufficient periods are proved. Their minimality does not follow
from those sufficient bounds, and a period in packet index is a degree shift
of $NP$.

\item \textbf{Exact affine Mahler-kernel valuations.}
For the kernel of Section~\ref{subsec:polyphase-exposition}, determine exact
values of
\[
v_p\bigl(\Theta_{M,s}(r,k)\bigr),\qquad
\Theta_{M,s}(r,k)=[z^r](1+z)^s((1+z)^M-1)^k,
\]
beyond the explicit special families. The minimum valuation of the summands
gives a lower bound. Equality requires a separate argument excluding
cancellation.

\item \textbf{Exact classical dyadic tail values.}
The uniform truncation exponent $4M+3$ and its optimality are proved.
Individual term weights are also known. The valuation of the whole omitted
tail at an arbitrary even row may be higher, so an exact formula for all rows
requires more than taking the minimum of those term weights. Empty tails
must remain exactly zero.
\end{enumerate}

These questions refine existing estimates and exact cases. They do not reopen
the external conjecture or the sharp uniform classical dyadic bound.

\subsection{Historical material and interface work}
\label{subsec:historical-scope}

The original status-corrected Paper IX survives as a 28-page PDF. Its recovered
formalization supports the framework, Euler coordinates, Mahler differences,
unit coordinates, recurrence packets, and CRT results retained in this
manuscript. The original TeX, original figure-generation files, and complete
historical experiment package were not recovered with that PDF. Their older
figure and test counts are therefore historical records, not newly reproduced
computations in this consolidation.

A separate reconstructed Annals text survives as a retrieved textual copy.
It includes finite-field mode and subgroup inversion, Fourier invariance,
Carmichael and matrix transfer, Amice pullbacks, depleted analytic Euler
branches, and extensive polyphase formulas. The sharp classical dyadic kernel
and truncation claims from that text have since been connected to checked
Lean proofs. The finite polyphase formulas used for explanation here have
short explicit algebraic derivations. The remainder of that larger historical
text has not acquired blanket formal coverage merely by being preserved.
The source archive retains the text with its provenance, and the statement
map distinguishes these historical claims from the certified claims in the
present manuscript.

Remaining recovery work includes the precise older odd-prime edge identities,
their stated ranges, and the binary digit-sum form of the checked weights.
Recovering a formula, checking its prose proof, and locating an exact Lean
declaration are separate steps. Missing historical sources should remain
identified instead of being reconstructed from a remembered result.

There is also a narrower interface task. The new regularity theorem applies to
arbitrary ambient $\Qp$ polynomial coefficients and identifies the analytic
germ of every continuous ambient-field root selection under its hypotheses.
Several older rectangular charts were defined using integer coefficients,
integer centers, or subtype-valued maps. Comparing each such historical
interface explicitly with the new ambient formulation remains to be done.
It does not invalidate the proved analytic implicit-function theorem, and it
is not counted as a general singular-persistence result.

\section{Formal evidence and reproduction}
\label{sec:verification}

The mathematical results in this manuscript have several distinct evidence
layers. The Lean source states the quantified theorem and its hypotheses. A
compiler report records acceptance of that source by a specified Lean and
Mathlib environment. An axiom audit identifies the axioms used by the resulting
named declarations. Statement contracts and mathematical controls check
selected interfaces and boundary cases. Finally, source correspondence reviews
compare the formal objects with the intended mathematical statement. These
layers answer different questions; a passing numerical example alone cannot
supply a universal theorem or establish that a formalization matches an
external conjecture.

This consolidation changes the exposition and the evidence index. It does not
change the mathematical Lean sources and does not claim a new compilation of
every preserved proof. Instead, the accompanying index checks that the source
files and retained logs agree with the exact hashes in their earlier successful
reports. The most recent compilation is the local analytic regularity release.
Its polynomial scope must be distinguished from the earlier arithmetic and
native obstruction compilations.

\subsection{Verification frontiers}

Table~\ref{tab:verification-frontiers} gives the recorded scope of each major
verification frontier. An \emph{inventoried module} is a source file in that
frontier's formalization list. Older lists sometimes include files of examples
or positive controls in this group. A \emph{check} is an entry in the report's
run list, including expected rejection of invalid claims. An \emph{axiom
closure} is the recorded axiom dependency set of an explicitly named
declaration; definitions and abbreviations contribute to this count as well
as theorems. None of these columns is a count of distinct mathematical results
across the paper.

\begin{table}[htbp]
\centering
\small
\begin{tabular}{@{}lrrr@{}}
\toprule
Preserved frontier & Modules & Checks & Axiom closures\\
\midrule
Original completion baseline & 79 & $79+37+3$ & 1000\\
Dyadic strengthening and isometry & 22 & 38 & 246\\
Standalone finite affine obstruction & 2 & 4 & 17\\
Shared correction quotient & 6 & 13 & 60\\
Polynomial one-step lifting & 9 & 18 & 91\\
Exact polynomial lift counts & 9 & 17 & 106\\
Smooth infinite towers & 14 & 22 & 127\\
Coefficient deformation & 17 & 26 & 159\\
Continuous parameter charts & 25 & 40 & 244\\
Weighted lifting & 30 & 50 & 304\\
First-order root motion & 35 & 58 & 326\\
Prescribed parameter obstruction & 39 & 65 & 375\\
Prescribed parameter towers & 45 & 74 & 408\\
Joint coefficient and parameter motion & 49 & 82 & 460\\
Local analytic regularity & 55 & 91 & 515\\
\bottomrule
\end{tabular}
\caption{Recorded verification scopes. Rows overlap extensively through
inherited sources and must not be added. The baseline separates 79 module
checks, 37 control checks, and three additional recorded runs; its nine
dependency revision comparisons are recorded separately.}
\label{tab:verification-frontiers}
\end{table}

The baseline inventory contains 836 theorem commands, 162 definitions and two
abbreviations, matching its 1000 audited names. The latest polynomial inventory
contains 322 theorem commands, 76 definitions and seven abbreviations in its
55 formalization modules. Its eight positive-control support modules and
13 statement-contract modules bring the full audit to 415 theorem commands,
92 definitions and eight abbreviations, hence 515 closures. The 91 checks also
include 12 negative controls. Six new regularity modules contribute 30 theorem
commands and eight definitions; all 49 inherited formalization modules are
unchanged in that release.

The arithmetic strengthening inventory exposes a small historical metadata
error. Its summary fields still say 179 theorems and 12 definitions, while its
actual module entries list 215 theorem commands, ten lemma commands,
20 definitions and one abbreviation. These 246 named entries exactly match
the preserved axiom audit. The consolidation uses the entries and audit, and
records the stale summary fields explicitly. The original inventory remains
unchanged as historical evidence. This is a correction to the reporting
metadata, not a newly discovered proof or a modification of a checked theorem.

\subsection{Claim-to-source correspondence}

The source references attached to the mathematical sections identify the
principal declarations for each result. Their complete names, exact source
hashes, roles and verification frontiers are collected in
\texttt{consolidation/audit/DECLARATIONS.csv}. The accompanying
\texttt{consolidation/audit/EVIDENCE\_INDEX.json} records the corresponding inventories,
reports and retained logs. In particular, the following interfaces distinguish
the principal completed results from generic auxiliary lemmas:
\begin{itemize}
\item The external conjecture is assembled by
\lean{NishibuchiCompletion.conjecture_3_9} in the original completion
baseline. Its coefficient identification, integrality, parameter split and
explicit threshold are part of that proof chain.
\item Exact $N=4$ packet periods are supplied by
\lean{DyadicFourLeastPeriods.least_packet_period}; the isometric extension
is supplied by \lean{DyadicFourIsometry.existsUnique_isometry}. The
normalization of the original coefficients is an explicit part of these
interfaces.
\item The classical sharp edge tower is supplied by
\lean{DyadicSharpTower.dyadic_supercongruence_sharp}, together with its
first omitted error valuation. Its optimality concerns the stated uniform
truncation bound.
\item Prescribed parameter lifting uses the actual residual and actual
partial Jacobians in
\lean{PolynomialPrescribedObstruction.prescribed_lift_iff_obstruction_zero}.
The quotient criterion and exact lift counts are not assumptions of the
later tower results.
\item Joint deformation is supplied by
\lean{PolynomialJointDeformation.exists_unique_joint_deformation}.
Local regularity is supplied by
\lean{PolynomialAnalyticImplicit.exists_analytic_implicit}, with the
actual polynomial derivative and the derivative formula $-J^{-1}K$.
\item \lean{PolynomialAnalyticSelection.continuous_selection_regularity}
transfers the analytic germ and strict derivative to any continuous local
root selection satisfying the stated root and base-point conditions.
\end{itemize}
Here arithmetic declaration names have namespace prefix
\texttt{KomorebiPaperIX.}; polynomial and native declarations have prefix
\texttt{Komorebi.}. The section maps and CSV retain the full names.

The declaration index covers every explicitly named theorem, lemma,
definition and abbreviation in the 14 release-frontier inventories, together
with the earlier standalone affine report. Byte-identical copies of a source
file are collapsed, while all their locations remain recorded. If the same
qualified name occurs in different source versions, the versions remain
separate and are marked. Anonymous examples, compiler-generated declarations,
failed development attempts, and the contents of earlier nested historical
snapshots are not counted as named results in this index. Negative-control
source files are checked for integrity but are not recorded as proved
mathematical assertions. Consequently, the number of CSV rows must not be
quoted as a global unique theorem count.

\subsection{External statement and library dependencies}

The original conjecture correspondence was reviewed against Nishibuchi's
arXiv version \texttt{2605.10677v2}, dated 1 June 2026. The retained source
record identifies the official version, the full title, the PDF checksum and
the extracted-text checksum. The detailed review compares Definition~2.1,
Notation~3.2 and Conjecture~3.9 with the actual reciprocal EGF coefficients,
the factor $j!$, the residue range, the degree displacement, and the threshold
in the formal capstone. Nishibuchi's conjecture is the conclusion to be
proved; it is not an axiom imported into Lean.

The full external paper is not redistributed in this source release.
Consolidation preserves the earlier source correspondence audit and its
version record. A fresh request to the official arXiv endpoints was blocked
by the retrieval tool, so this edition does not claim a new inspection of the
primary paper. This access failure affects fresh reauthentication of the
external wording, not the recorded compiler acceptance of the unchanged
formal sources. No broad literature or priority claim is made.

The mathematical environment is Lean \texttt{4.34.0-rc2}, with Mathlib revision
\nolinkurl{aeb72f1e9981b6d35b426b16dca7759be3e7ba06}.
The allowed axiom dependencies are
\texttt{propext}, \texttt{Classical.choice} and \texttt{Quot.sound}.
The reports do not describe a rebuild of Mathlib from source. The local
regularity audit separately records the applicable Mathlib analytic inverse
and strict derivative interfaces and their hypotheses. In particular, the
field-valued analysis is performed over a complete nontrivially normed field;
its integral specialization embeds the data into $\Qp$. It does not use
$\Zp$ as a vector space over $\Qp$.

\subsection{Reproduction and limits of the consolidation check}

The archive contains the mathematical sources, verifier scripts, controls,
contracts, prior logs and exact inventories. The pinned Lean and Mathlib kit
is supplied separately. From the consolidated archive root, the latest
polynomial gate can be reproduced with a new output directory:
\begin{verbatim}
python3 replay_paper9_regularity.py \
  --kit /absolute/path/to/pinned-kit \
  --output /absolute/path/to/new-polynomial-run
\end{verbatim}
The preserved arithmetic, native and shared-correction chain has a separate
entry point:
\begin{verbatim}
python3 replay_shared_correction.py --all \
  --kit /absolute/path/to/pinned-kit \
  --output /absolute/path/to/new-arithmetic-run
\end{verbatim}
These commands are documented reproduction entry points, not additional
runs performed during manuscript consolidation. The second command rebuilds
the preserved baseline before the arithmetic continuation. Each output path
must be new. The reports retain the original full commands, runtime
fingerprints and output hashes. The executable-path compatibility adapter
used in the recorded environment restores runtime path discovery; it does
not alter Lean's mathematical kernel or the source statements.

For the narrower evidence reconciliation carried out for this edition, use
\begin{verbatim}
python3 consolidation/tools/evidence_index.py \
  --source . --output /absolute/path/to/new-evidence-index
\end{verbatim}
This check compared 644 source-file occurrences with their recorded hashes
and 713 retained log occurrences with their recorded hashes. Four additional
logs from the older standalone affine report were fingerprinted, but that
legacy report does not contain prior log hashes; those four are not promoted
to hash comparisons. Every named declaration in the selected inventories
matches the appropriate recorded standard-axiom dependency set. These are
integrity and provenance checks, not another Lean compilation.

Expected negative-control failures are retained as evidence that specific
invalid strengthenings are rejected. Examples include an unjustified extra
precision digit, omitted parameter terms, a wrong derivative sign, and
uniqueness without the appropriate hypotheses. Development failures and
interrupted runs remain historical records; only the successful designated
release reports support the verification table. A compilation failure is
never counted as a proved theorem merely because its source resembles a
successful one.

Finally, the proof of the original stated conjecture does not resolve every
extension mentioned in this manuscript. The open questions concerning
rationally rank-deficient singular points, presentation invariance, global
or completed descent, and further exact arithmetic have no certificate
supplied by these counts. Their status is determined by their statements and
missing implications, not by the size of the archive.


\begin{thebibliography}{9}
\bibitem{Nishibuchi}
Yuta Nishibuchi,
\emph{$p$-adic Congruences of Generalized Euler Numbers and Relations to Even
Zeta Values}, arXiv:2605.10677v2, 1 June 2026.
\url{https://arxiv.org/abs/2605.10677v2}.
The comparison to Conjecture 3.9 uses the source review and fixed-version
record preserved with this edition.

\bibitem{Mathlib}
The Mathlib Community,
\emph{Mathlib}, source revision
\nolinkurl{aeb72f1e9981b6d35b426b16dca7759be3e7ba06},
used with Lean 4.34.0-rc2.
In particular, the polynomial, strict inverse-function and analytic
inverse-function developments cited in Section~\ref{sec:regularity}.
The exact local source paths and dependency records accompany the release.

\bibitem{PaperIXArchive}
Komorebi P.,
\emph{Projector-Resolved Seed Calculus IX: original manuscript, formal
supplements and cumulative continuation sources}, preserved source collection,
consolidated 8 September 2026.
The original 28-page status-corrected manuscript and all subsequent release
records remain in the accompanying archive.
\end{thebibliography}
\end{document}
